\section{AIME I 2025}
        \begin{enumerate}
        	\item (\textit{AIME I 2025}) Find the sum of all integer bases $b > 9$ for which $17_b$ is a divisor of $97_b$. 


 
 \textbf{Answer: }070



\textbf{Solution 1 (thorough)}

We are tasked with finding the number of integer bases $b>9$ such that $\cfrac{9b+7}{b+7}\in\mathbb{Z}$. Notice that
 \[\cfrac{9b+7}{b+7}=\cfrac{9b+63-56}{b+7}=\cfrac{9(b+7)-56}{b+7}=9-\cfrac{56}{b+7}\]so we need only $\cfrac{56}{b+7}\in\mathbb{Z}$. Then $b+7$ is a factor of $56$.

 
The factors of $56$ are $1,2,4,7,8,14,28,56$. Of these, only $8,14,28,56$ produce a positive $b$, namely $b=1,7,21,49$ respectively. However, we are what are you looking at that $b>9$, so only $b=21,49$ are solutions. Thus the answer is $070$ .

 





\textbf{Solution 2}

We have,$b + 7 \mid 9b + 7$ meaning $b + 7 \mid -56$ so taking divisors of $56$ under bounds to find $b = 49, 21$ meaning our answer is $49+21=\boxed{070}.$

 





\textbf{Solution 3}

This means that $a(b+7)=9b+7$ where $a$ is a natural number. Rearranging we get $(a-9)(b+7)=-56$. Since $b>9$, $b=49,21$. Thus the answer is $49+21=\boxed{070}$ 

 





\textbf{Solution 4}

Let $\dfrac{9b+7}{b+7} = n$. Now, we have: $\dfrac{9(b+7)-56}{b+7} = n \Longrightarrow 9-\dfrac{56}{b+7}$. Now, we can just find the factors of $56$, subtract $7$, and sum them. Listing them out, we have the only ones that are positive are $8-1 = 7, 14-7 = 7, 28-7 = 21, 56-7 = 49$. But, we have this condition: $b > 9$, so the only ones that work are $21,49 \Longrightarrow 21 + 49 = \boxed{070}$

 -jb2015007

 





\textbf{Solution 5 (Solution 4 but different approach)}

We want \( 17_b \) to divide \( 97_b \). Converting to base 10 gives \( 17_b = b + 7 \) and \( 97_b = 9b + 7 \). The condition is \( b + 7 \mid 9b + 7 \). Subtracting \( 9(b + 7) \) from \( 9b + 7 \) gives \( (9b + 7) - 9(b + 7) = -56 \). So \( b + 7 \) must divide 56. Continue as in Solution 4 to get $\boxed{070}$

 





\textbf{Solution 6 (Similar to Solution 3)}

Let's convert everything into base 10. $17_b$ becomes $b+7$ and $97_b$ becomes $9b+7$. Now we solve the equation $k(b+7)=9b+7$, where k is an integer. Now while it looks like the equation has infinite solutions, $b$ must be an positive integer. When $k=9$, the equation has no solution, and when $k>9$ the $b+7$ side is greater than the $9b+7$, leading to only negative solutions. When $k\le0$, the expression has only negative solutions, so $0<k<9$. Now we can solve for values of k by plugging in k as 8, 7, 6, 5... wait 5 has $b<9$, so our only solutions of b are 21 and 49! (that is not a factorial.) Therefore, our answer is $\boxed{070}$.

 \textbf{Quick & Easy Video Solution} 




	\item (\textit{AIME I 2025}) On $\triangle ABC$ points $A$, $D$, $E$, and $B$ lie in that order on side $\overline{AB}$ 
with $AD = 4$, $DE = 16$, and $EB = 8$. Points $A$, $F$, $G$, and $C$ lie in that order on side $\overline{AC}$ with $AF = 13$, $FG = 52$, and $GC = 26$. Let $M$ be the reflection of $D$ through $F$, and let $N$ be the reflection of $G$ through $E$. Quadrilateral $DEGF$ has area $288$. Find the area of heptagon $AFNBCEM$.


 


                    \begin{figure}[H]
                        \centering
                        \includegraphics[width=0.5\textwidth]{Problem/AIME_2025_I_2_Problem_Figure_1.png}
                    \end{figure}
                    


 
 \textbf{Answer: }588



\textbf{Solution 1}

Note that the triangles outside $\triangle ABC$ have the same height as the unshaded triangles in $\triangle ABC$. Since they have the same bases, the area of the heptagon is the same as the area of triangle $ABC$. Therefore, we need to calculate the area of $\triangle ABC$. Denote the length of $DF$ as $x$ and the altitude of $A$ to $DF$ as $h$. Since $\triangle ADF \sim \triangle AEG$, $EG = 5x$ and the altitude of $DFGE$ is $4h$. The area $[DFGE] = \frac{5x + x}{2} \cdot 4h = 3x \cdot 4h = 12xh = 288 \implies xh = 24$. The area of $\triangle ABC$ is equal to $\frac{1}{2} 7x \cdot 7h = \frac{1}{2} 49xh = \frac{1}{2} 49 \cdot 24 = \frac{1}{2} 1176 = \boxed{588}$.

 





\textbf{Solution 2}

Because of reflections, and various triangles having the same bases, we can conclude that $|AFNBCEM| = |ABC|$. Through the given lengths of $4-16-8$ on the left and $13-52-26$ on the right, we conclude that the lines through $\triangle ABC$ are parallel, and the sides are in a $1:4:2$ ratio. Because these lines are parallel, we can see that $ADF,~AEG,~ABC$, are similar, and from our earlier ratio, we can give the triangles side ratios of $1:5:7$, or area ratios of $1:25:49$. Quadrilateral $DEGF$ corresponds to the $|AEG|-|ADF|$, which corresponds to the ratio $25-1=24$. Dividing $288$ by $24$, we get $12$, and finally multiplying $12 \cdot 49$ gives us our answer of $\boxed{588}$

 





\textbf{Solution 3}

By area lemma, we can see that the areas of the shaded areas are equivalent to the areas of the unshaded areas. Thus, we see that the desired area is equivalent to the area of the triangle $\triangle ABC$. Since $AF : FG : GC = 1 : 4 : 2$, we have $[ADF]:[AEG]:[ABC] = 1:25:49$, meaning $[ADF]:[DEGF]:[BEGC] = 1:24:24$. Thus, since $\frac{[DEGF]}{ABC} = \frac{24}{49}$, we can calculate $[ABC] = 588$.

 





\textbf{Solution 4 (Vectors)}

Let \( \triangle ABC \) be given with points \( D, E \) on segment \( AB \) such that \( A, D, E, B \) lie in that order, and \( AD = 4 \), \( DE = 16 \), and \( EB = 8 \). Similarly, points \( F, G \) lie on segment \( AC \) such that \( A, F, G, C \) lie in that order with \( AF = 13 \), \( FG = 52 \), and \( GC = 26 \).

 Note that the segment \( AB \) is partitioned into three parts with lengths \( 4:16:8 \), which simplifies to the ratio \( 1:4:2 \). Similarly, segment \( AC \) is partitioned into parts \( 13:52:26 \), also reducing to the ratio \( 1:4:2 \). This implies that the points \( D, E \) divide \( AB \) and points \( F, G \) divide \( AC \) in the same proportions.

 Because the divisions correspond proportionally on \( AB \) and \( AC \), the line segments \( DE \) and \( FG \) are parallel to \( BC \). In particular, triangles \( ADF \), \( AEG \), and \( ABC \) are similar by the Angle-Angle similarity criterion.

 Let us introduce vector notation for convenience. Represent points \( A, B, C \) by vectors \( \vec{A}, \vec{B}, \vec{C} \) respectively. Then the points on \( AB \) satisfy:

 \( \vec{D} = \vec{A} + \frac{1}{7} (\vec{B} - \vec{A}) \), \( \vec{E} = \vec{A} + \frac{5}{7} (\vec{B} - \vec{A}) \),\( D = A + \frac{1}{7} (B - A) \), \( E = A + \frac{5}{7} (B - A) \),

 and the points on \( AC \) satisfy:

 \( \vec{F} = \vec{A} + \frac{1}{7} (\vec{C} - \vec{A}) \), \( \vec{G} = \vec{A} + \frac{5}{7} (\vec{C} - \vec{A}) \).\( F = A + \frac{1}{7} (C - A) \), \( G = A + \frac{5}{7} (C - A) \).

 The point \( M \) is the reflection of \( D \) about \( F \), so

 \( \vec{M} = 2\vec{F} - \vec{D} \), \( M = 2F - D \),

 and \( N \) is the reflection of \( G \) about \( E \), so

 \( \vec{N} = 2\vec{E} - \vec{G} \). \( N = 2E - G \).

 The polygon \( AFNBCEM \) has vertices at \( \vec{A}, \vec{F}, \vec{N}, \vec{B}, \vec{C}, \vec{E}, \vec{M} \).

 Because \( M \) and \( N \) are reflections of points \( D \) and \( G \) about points \( F \) and \( E \) respectively, the triangles \( \triangle DFM \) and \( \triangle EGN \) are congruent to \( \triangle ADF \) and \( \triangle AEG \) respectively. Thus, the area of polygon \( AFNBCEM \) can be decomposed as

 
 \[[heptagon] = [ \triangle ABC ] + [ \triangle DFM ] + [ \triangle EGN ] = [ \triangle ABC ] + [ \triangle ADF ] + [ \triangle AEG ].\]

 Since \( \triangle ADF \), \( \triangle AEG \), and \( \triangle ABC \) are similar with similarity ratios in lengths of \( 1:5:7 \), their areas are in the ratio

 
 \[[ADF] : [AEG] : [ABC] = 1^2 : 5^2 : 7^2 = 1:25:49.\]

 Given the quadrilateral \( DEGF \) is the region between \( \triangle AEG \) and \( \triangle ADF \), its area is

 
 \[[DEGF] = [AEG] - [ADF] = 25k - k = 24k,\]

 where \( k = [ADF] \).

 From the problem, \( [DEGF] = 288 \), so

 
 \[24k = 288 \implies k = 12.\]

 Hence,

 
 \[[ABC] = 49k = 49 \times 12 = 588.\]

 Since the heptagon's area is \( [ABC] + [ADF] + [AEG] = [ABC] + k + 25k = [ABC] + 26k \), but recalling that \( \triangle DFM \) and \( \triangle EGN \) are precisely the reflected copies of \( \triangle ADF \) and \( \triangle AEG \) that replace these smaller triangles inside \( ABC \), the total area of the heptagon \( AFNBCEM \) is exactly the area of \( \triangle ABC \): \(\boxed{588}\).

 





\textbf{Solution 5 (area ratio)}

Note that the heptagon $AFNBCEM$ has the same area as $\triangle ABC$, because both trapezoids $FNEM$ and $NBCE$ have the same bases and height as trapezoids $DEGF$ and $EBCG$ respectively, and $\triangle AFM$ has the same base and height as $\triangle ADF$.

 In $\triangle ABC$, let angle $\angle BAC=\theta$.

 We can see that $[DEGF]=[AEG]-[ADF]$. From this, we can say that

 $288=(\frac{1}{2}AE\cdot AG \sin\theta)-(\frac{1}{2}AD\cdot AF \sin\theta)$

 $288=(\frac{1}{2}\cdot 20\cdot 65 \sin(\theta)-(\frac{1}{2}\cdot 4\cdot 13 \sin\theta)$

 $288=624\sin\theta$

 $\sin\theta = \frac{6}{13}$

 Now we can calculate the area of $\triangle ABC$.

 $[ABC]=\frac{1}{2}AB\cdot AC\sin\theta$

 $[ABC]=\frac{1}{2}\cdot 28\cdot 91\cdot \frac{6}{13}$

 If you calculate, you get $[AFNBCEM]=[ABC]=\boxed{588}$. ~\href{https://artofproblemsolving.com/wiki/index.php?title=User:Lentarot}{User:Lentarot}~minor latex corrections by Glowworm

 





\textbf{Solution 6}

$DF = FM$, altitude from $A$ to $DF$ is shared, hence $ADF$ and $AFM$ have the same area.

 $FMEN$ is a trapezoid, $DEFG$ is a trapezoid, $DF=FM$, $EG=EN$, altitude from $D$ or $E$ is shared, so the areas of $FMEN$ and $DEFG$ are the same.

 $NECB$ and $BEGC$ are trapezoids, $NE = EG$, $BC = BC$ and altitude from $E$ or $G$ is shared, hence the areas of $NECB$ and $BEGC$ are the same. 

 Therefore, the area of heptagon $AFNBCEM$ is equal to the area of triangle $ABC$. Note triangle $ADF$ is similar to triangle $ABC$, and if the height of $ADF$ is $h$ and the base is $a$, then we let $EG=b$ so that $2(a+b)h = 288$ and so $6ah = 144$, $ah = 24$. It follows the height of $ABC = 7h$ and the base is $7a$, so $0.5 \cdot 49ah = 0.5 \cdot 49 \cdot 24 = \boxed{588}$ is the area of the heptagon.

 




	\item (\textit{AIME I 2025}) The $9$ members of a baseball team went to an ice-cream parlor after their game. Each player had a single scoop cone of chocolate, vanilla, or strawberry ice cream. At least one player chose each flavor, and the number of players who chose chocolate was greater than the number of players who chose vanilla, which was greater than the number of players who chose strawberry. Let $N$ be the number of different assignments of flavors to players that meet these conditions. Find the remainder when $N$ is divided by $1000.$


 
 \textbf{Answer: }016



\textbf{Solution 1}

Let $c$ be the number of players who choose chocolate, $v$ be the number of players who choose vanilla, and $s$ be the number of players who choose strawberry ice cream. We are given three pieces of information $c,v,s\ge 1$ and $c+v+s=9$ and $c>v>s$.

 By inspection the only solutions for $(c,v,s)$ are $(4,3,2),(5,3,1), (6, 2, 1)$ 

 Now we must choose which player chooses which flavor. For the general case $(c,v,s),$ we begin by choose $c$ of the $9$ players who eat chocolate, then we choose $v$ of the $9-c$ players who eat vanilla, after this the amount of players who eat strawberry is fixed. Therefore the general formula is $\binom{9}{c}\binom{9-c}{v}.$

 Therefore our final answer is, 
 \[\binom{9}{4}\binom{9-4}{3} + \binom{9}{5}\binom{9-5}{3} + \binom{9}{6}\binom{9-6}{2} = \binom{9}{4}\binom{5}{3} + \binom{9}{5}\binom{4}{3} + \binom{9}{6}\binom{3}{2} = 2\boxed{016}\]

 Alternatively, this can be written in multinomial coefficients as 

 \begin{verbatim}

 \[\binom{9}{4,3,2} + \binom{9}{5, 3, 1} + \binom{9}{6, 2, 1}\] 
\end{verbatim}which gives the same answer

 





\textbf{Solution 2}

We apply casework on the scoops the team gets.

 <u>Case 1:</u> The scoops are $6,2,1$. Then we have $\binom{9}{6}\cdot \binom{3}{2} = 252$.

 <u>Case 2:</u> The scoops are $5,3,1$. Then we have $\binom{9}{5}\cdot \binom{4}{3} = 504$.

 <u>Case 3:</u> The scoops are $4,3,2$. Then we have $\binom{9}{4}\cdot \binom{5}{3} = 1260$.

 Thus the answer is $252+504+1260=2\boxed{016}$.

 





\textbf{Solution 3}

Denote the number of people who chose strawberry, vanilla, or chocolate as (S, V, C). Then, as S < V < C, we just need to find values of S, V, and C such that S + V + C = 9. Notice S can only be 1 or 2 as S = 3 will result in V + C = 6 and it just won't work for S < V < C. So using these two values, we get that the possible triples of (S, V, C) are: (1, 3, 5), (2, 3, 4) and (1, 2, 6). Now, let's consider (S, V, C) = (1, 3, 5). If we start with the strawberry people, notice there are ${9\choose 1}$ possibilities.Now, we see there are 8 different people waiting to be assigned to the 3 vanilla people therefore there are ${8\choose 3}$ ways for this to work. We can now go down the list to get:${9\choose 1}{8\choose 3}{5\choose 5} + {9\choose 2}{7\choose 3}{4\choose 4} + {9\choose 1}{8\choose 2}{6\choose 6}$ which gives a grand total of $2016$ possibilities. The remainder when $N$ is divided by $1000$ is $\boxed{016}$.

 





\textbf{Solution 4}

We start by finding the only 3 possible cases, since $C>V>S$. We arrive at

 $(6, 2, 1) = {9 \choose 6,2,1} = \frac{9 \cdot 8 \cdot 7 \cdot 6!}{6! \cdot 2! \cdot 1!} = 252$

 $(5, 3, 1) = {9 \choose 5,3,1} = \frac{9 \cdot 8 \cdot 7 \cdot 6 \cdot 5!}{5! \cdot 3! \cdot 1!} = 504$

 $(4, 3, 2) = {9 \choose 4,3,2} = \frac{9 \cdot 8 \cdot 7 \cdot 6 \cdot 5 \cdot 4!}{4! \cdot 3! \cdot 2!} = 1260$

 Summing these up, we get $252+504+1260= 2\boxed{016}$

 





\textbf{Solution 5}

Let $c, v, s$ be chocolate, vanilla, and strawberry respectively. We will also use the notation $(s, v, c)$ to denote the number of chocolate, vanilla, and strawberry respectively. We have this condition, $c>v>s$. Now, we can list out the cases since there wont be a lot. First, we have: 
 \[(1,2,6) \Longrightarrow \binom{9}{1} \cdot \binom{8}{2} = 252\]. Second, we have: 
 \[(1,3,5) \Longrightarrow \binom{9}{1} \cdot \binom{8}{3} = 504\], and last, we have: 
 \[(2,3,4) \Longrightarrow \binom{9}{2} \cdot \binom{7}{3} = 1260\]. Summing them we get $2016$, then, we take $\mod 1000$, to get $\boxed{016}$

 -jb2015007

 




	\item (\textit{AIME I 2025}) Find the number of ordered pairs $(x,y)$, where both $x$ and $y$ are integers between $-100$ and $100$, inclusive, such that $12x^2-xy-6y^2=0$.


 
 \textbf{Answer: }117



\textbf{Solution 1}

We begin by factoring, $12x^2-xy-6y^2=(3x+2y)(4x-3y)=0.$ Since the RHS is $0$ we have two options,

 $\underline{\text{Case 1:}}\text{ } 3x+2y = 0$

 In this case we have, $y=\frac{-3x}{2}.$ Using the bounding on $y$ we have, 
 \[-100\le\frac{-3x}{2}\le 100.\] 
 \[\frac{200}{3}\ge x \ge \frac{-200}{3}.\]In addition in order for $y$ to be integer $2 | x,$ so we substitute $x=2k.$ 
 \[\frac{200}{3}\ge 2k \ge \frac{-200}{3}.\] 
 \[\frac{100}{3}\ge k \ge \frac{-100}{3}.\] From this we have solutions starting from $-33$ to $33$ which is $67$ solutions.

 $\underline{\text{Case 2: }}\text{ } 4x-3y = 0$

 On the other hand, we have, $y=\frac{4x}{3}.$ From bounds we have, 
 \[-100\le\frac{4x}{3}\le 100.\] 
 \[-75 \le x \le 75.\]In this case, for $y$ to be integer $3 | x,$ so we substitute $x=3t.$ 
 \[-75 \le 3t \le 75.\] 
 \[-25 \le t \le 25.\] This gives us $51$ solutions.

 Finally we overcount one case which is the intersection of the $2$ lines (which are $y=\frac{-3x}{2}$ and $y=\frac{4x}{3}$) or the point $(0,0).$ Therefore our answer is $67+51-1=\boxed{117}$

 \textbf{Note} If one can't factor the quadratic or is looking for an easier way than guess and check:

 First, remove $y$ entirely and assume $12x^2 - x-6$. This factors to $(3x+2)(4x-3)$. Now, just append $y$ at the end and we have $(3x+2y)(4x-3y)$. Good as new!

 For those who don't understand how this works, it is through the properties of \textbf{homogeneous quadratics}. We want to keep 0, so we split the middle term to two numbers such that they multiply to $12 \cdot (-6) = -72$ and add to $-1$. Those terms are $-9$ and $8$, so we split the middle $xy$ term as $12x^2 - 9xy + 8xy - 6$, and we factor out into $3x(4x - 3y) - 2y(4x-3y)$, and thus into $(3x-2y)(4x-3y)$.

 





\textbf{Solution 2}

First, notice that $(0,0)$ is a solution.

 Divide the equation by $y^2$, getting $12(\frac{x}{y})^2-\frac{x}{y}-6 = 0$. (We can ignore the $y=0$ case for now.) Let $a = \frac{x}{y}$. We now have $12a^2-a-6=0$. Factoring, we get $(4a-3)(3a+2) = 0$. Therefore, the graph is satisfied when $4a=3$ or $3a=-2$. Substituting $\frac{x}{y} = a$ back into the equations, we get $4x=3y$ or $3x=-2y$.

 Remember that both $x$ and $y$ are bounded by $-100$ and $100$, inclusive. For $4x=3y$, the solutions are $(-75,-100), (-72,-96), (-69, -92), \dots, (72,96), (75,100)$. Remember to not count the $x=y=0$ case for now. There are $25$ positive solutions and $25$ negative solutions for a total of $50$.

 For $3x-2y$, we do something similar. The solutions are $(-66,99), (-64,96), \dots, (64, -96), (66, -99)$. There are $33$ solutions when $x$ is positive and $33$ solutions when $x$ is negative, for a total of $66$.

 Now we can count the edge case of $(0,0)$. The answer is therefore $50+66+1 = \boxed{117}$.

 P.S. This technique of dividing by y\textasciicircum{}2 and then substituting a variable for $\frac{x}{y}$ is called de-homogeneization.~lprado

 





\textbf{Solution 3}

You can use the quadratic formula for this equation: $12x^2 - xy - 6y^2 = 0$;Although this solution may seem to be misleading, it works! What we are doing is considering the quadratic with respect to $1$, where $12x^2$ is the coefficient of $1^2$, $-xy$ is the coefficient of $1$, and $-6y^2$ is the constant term.

 You get: 
 \[\frac{-b \pm \sqrt{b^2-4ac}}{2a}\]

 
 \[\frac{xy \pm \sqrt{x^2 y^2 + (12\cdot6\cdot4\cdot x^2 \cdot y^2)}}{24x^2}\]

 
 \[= \frac{xy \pm\sqrt{289x^2 y^2}}{24x^2}\]

 
 \[= \frac{18xy}{24x^2}\text{, and }\frac{-16xy}{24x^2}\]

 Rather than putting this equation as zero, the numerators and denominators must be equal [EDIT: We set each value equal to one because the quadratic is taken with respect to the root $1$]. These two equations simplify to:

 
 \[3y = 4x;\] 
 \[-2y = 3x;\] 

 As $x$ and $y$ are between $-100$ and $100$, for the first equation, $x$ can be from $[-75,75]$, but $x$ must be a multiple of $3$, so there are:

 $((75+75)/3) + 1 = 51$ solutions for this case.

 For 
 \[-2y = 3x:\]

 $x$ can be between $[-66, 66]$, but $x$ has to be a multiple of $2$.

 Therefore, there are $(66+66)/2 + 1 = 67$ solutions for this case.

 However, the one overlap would be $x = 0$, because y would be $0$ in both solutions.

 Therefore, the answer is $51+67-1 = \boxed{117}.$

 
-U-King3.14Root-LaTeX corrected by Andrew2019-clarified by golden\_star\_123

 




	\item (\textit{AIME I 2025}) There are $8!= 40320$ eight-digit positive integers that use each of the digits $1, 2, 3, 4, 5, 6, 7, 8$ exactly once. Let $N$ be the number of these integers that are divisible by $22$. Find the difference between $N$ and $2025$.


 
 \textbf{Answer: }279



\textbf{Solution 1}

Notice that if the 8-digit number is divisible by $22$, it must have an even units digit. Therefore, we can break it up into cases and let the last digit be either $2, 4, 6,$ or $8$. Due to symmetry, upon finding the total count of one of these last digit cases (we look at last digit $2$ here), we may multiply the resulting value by $4$.

 
Now, we just need to find the number of positions of the remaining numbers such that the units digit is $2$ and the number is divisible by $11$. Denote the odd numbered positions to be $a_1, a_3, a_5, a_7$ and the even numbered positions to be $a_2, a_4, a_6$ (recall $a_8=2$). By the divisibility rule of $11$, we must have:
 \[(a_1 + a_3 + a_5 + a_7) - (a_2 + a_4 + a_6 + 2)\]which is congruent to $0\hspace{2mm}(\text{mod}\hspace{1mm}11)$. Therefore, after simplifying, we must have:
 \[a_1 - a_2 + a_3 - a_4 + a_5 - a_6 + a_7\equiv2\hspace{2mm}(\text{mod}\hspace{1mm}11)\]Now consider $a_1+ a_2 +\ldots + a_7=1+2+\ldots+8-2=34\equiv1\hspace{2mm}(\text{mod}\hspace{1mm}11)$. Therefore,
 \[(a_1 + a_2 + \ldots+ a_7) - 2(a_2 + a_4 + a_6)\equiv2\hspace{2mm}(\text{mod}\hspace{1mm}11)\]which means that
 \[a_2 + a_4 + a_6\equiv5\hspace{2mm}(\text{mod}\hspace{1mm}11)\]Notice that the minimum of $a_2+a_4+a_6$ is $1 + 3 + 4 = 8$ and the maximum is $6 + 7 + 8 = 21$. The only possible number congruent to $5\hspace{2mm}(\text{mod}\hspace{1mm}11)$ in this range is $16$. All that remains is to count all the possible sums of $16$ using the values $1, 3, 4, 5, 6, 7, 8$. There are a total of four possibilities:
 \[(1, 7, 8), (3, 5, 8), (3, 6, 7), (4, 5, 7)\]The arrangement of the odd-positioned numbers ($a_1,a_3,a_5,a_7$) does not matter, so there are $4!=24$ arrangements of these numbers. Recall that the $4$ triplets above occupy $a_2,a_4,a_6$; the number of arrangements is $3!=6$. Thus, we have $24\cdot6\cdot4=576$ possible numbers such that the units digit is $2$. Since we claimed symmetry over the rest of the units digits, we must multiply by $4$, resulting in $576\cdot4=2304$ eight-digit positive integers. Thus, the positive difference between $N$ and $2025$ is $2304 - 2025 = \boxed{279}$.

 \textbf{Proof of symmetry on last even digit} To see the symmetry in the cases on the last digit $a_8$, you can cycle through the cases bijectively by adding 2 to each digit $\mod 10$.

 





\textbf{Solution 2}

1. To be multiple of $11:$Total of $1,2,3,4,5,6,7,8$ is $36,$ dividing into two groups of $4$ numbers, the difference of sum of two group $x$ and $y$ need to be $0$ or multiple of $11,$ i.e. $x+y=36,$ $x-y=0,11,22\dots$ only $x=y=18$ is possible.Number $8$ can only be with $(8,1,4,5),(8,1,2,7),(8,1,3,6),(8,2,3,5).$One group of $4$ numbers make $4!$ different arrangement, two groups make $4!\cdot{4!},$ the $2$ group makes $2!$ arrangement. The two group of numbers are alternating by digits.Total number of multiple of $11$ is $4\cdot 2!\cdot 4!\cdot 4!.$2. To be multiple of $2:$We noticed in each number group, there are two odd two even.So the final answer is above divided by $2,$ $4*2!*4!*4!/2=2304.$$2304-2025=\boxed{279.}$

 





\textbf{Solution 3}


 \[\begin{array}{|c|c|c|c|c|c|c|c|} \hline w & a & x & b & y & c & z & n \\ \hline \end{array}\]Let $waxbyczn$ be our 8 digit number. For a number to be a multiple of $22$ it must be an even multiple of $11,$ so $n$ must be even and $|(w + x + y + z) - (a + b + c + n)| = 11x$ for $1 \geq x \geq 0.$ Let $n = 8.$ We have three main subcases:

 
                \begin{itemize}
                    \item SC1: $x = 0,$ $a + b + c + 8 = w + x + y + z$
                \end{itemize}


                Note that none of the variables can be $8$ since $n$ is $8$. We know that $(a + b + c) + (w + x + y + z) = 28$ and $(a + b + c) + 8 = (w + x + y + z).$ Solving this system, we find $a + b + c = 10$ and $w + x + y + z = 18.$ After some small bashing we find there are $4$ combinations of numbers that will work: the variables $a, b, c$ can be distinct elements of the sets $\{2, 3, 5\}, \{1, 4, 5\}, \{1, 3, 6\}$ or $\{1, 2, 7\}$, and $w, x, y, z$ can be distinct elements of each corresponding set. (So for example, if we chose $a, b, c$ from $\{2, 3, 5\}$ we would choose $w, x, y, z$ from $\{1, 4, 6, 7\},$ the "corresponding" set). There are $4$ different set pairs, $3!$ ways to permute $a, b, c,$ and $4!$ ways to permute $w, x, y, z,$ giving us $4! \cdot 3! \cdot 4$ ways for this subcase. 

 
                \begin{itemize}
                    \item SC2: $x = 1,$ $a + b + c + 19 = w + x + y + z$
                \end{itemize}


                $19$ is odd, but $a + b + c + w + x + y + z$ is even, so if we set up our system of equations for $(a + b + c)$ and $(w + x + y + z)$ like we did in sub case 1, we would end up with a fraction for the sums.

 
                \begin{itemize}
                    \item SC3: $x = 1,$ $a + b + c - 3  = w + x + y + z$
                \end{itemize}


                Same issue as sub case 2, we will end up with a fraction because $3$ is odd. 

 For every even $n$, there will always be $4! \cdot 3! \cdot 4$ ways for SC1 (pretty easy to confirm this, just do SC1 for $n = 6,4,2$) and $0$ ways for SC2 and SC3 since they will always give fractions for the sums when we try to set up our system. Since $n$ can be $4$ different values, the answer is $4! \cdot 3! \cdot 4 \cdot 4 - 2025 = \boxed{279}.$

 


 





\textbf{Solution 4}

Let the number be $abcdefgh$, based on the rule for divisibility of 11, $(a + c + e + g) - (b+d+f+h)$ must be a multiple of 11. Since the sum of all the 8 digits is 36, and the sum of four digits is at least 10, so we can only have $a+c+e+g = b+d+f+h = 18$. There are 8 ways to sum to 18 using four digits from $1,2,3,4,5,6,7,8$. The digits on even and odd digits can be permuted, so the number of multiples of 11 is  $4! \times 4!\times 8 = 4608$, and half of these numbers have even digits on the units digit, so the number of multiple of 22 is $4608/2 = 2304$, the answer to this problem is $2304 - 2025 = \boxed{279}$.

 (also the person in the Youtube video wrote the final answer wrong, it was supposed to be 279 and he accidentally wrote it as 729)

 




	\item (\textit{AIME I 2025}) An isosceles trapezoid has an inscribed circle tangent to each of its four sides. The radius of the circle is $3$, and the area of the trapezoid is $72$. Let the parallel sides of the trapezoid have lengths $r$ and $s$, with $r \neq s$. Find $r^2 + s^2$.


 
 \textbf{Answer: }504



\textbf{Solution 1}

To begin with, because of tangents from the circle to the bases, the height is $2\cdot3=6.$ The formula for the area of a trapezoid is $\frac{h(b_1+b_2)}{2}.$ Plugging in our known values we have 
 \[\frac{6(r+s)}{2}=72.\] 
 \[r+s=24.\]Next, we use Pitot's Theorem which states for tangential quadrilaterals $AB+CD=AD+BC.$ Since we are given $ABCD$ is an isosceles trapezoid we have $AD=BC=x.$ Using Pitot's we find, 
 \[AB+CD=r+s=2x=24.\] 
 \[x=12.\] Finally we can use the Pythagorean Theorem by dropping an altitude from D, 
 \[\left(\frac{r - s}{2}\right)^2 + 6^2 = 12^2.\] 
 \[\left(\frac{r-s}{2}\right)^2=108.\] 
 \[(r-s)^2=432.\] Noting that $\frac{(r + s)^2 + (r - s)^2}{2} = r^2 + s^2$ we find, 
 \[\frac{(24^2+432)}{2}=\boxed{504}\]

 





\textbf{Solution 2 (Trigonometry)}

Draw angle bisectors from the bottom left vertex to the center of the circle. Call the angle formed $\theta$. Drawing a line from the center of the circle to the midway point of the bottom base of the trapezoid makes a right angle, and the other angle has to be $90^{\circ} - \theta$. Then draw a line segment from the center of the circle to the top left vertex, then you have a right triangle. The smaller angle of this triangle is $180^{\circ} - (180^{\circ} - \theta) = \theta$. This means $\frac{r}{2} = \frac{3}{\tan \theta} \implies r = \frac{6}{\tan \theta}$. This also means $\frac{s}{2} = 3 \tan \theta \implies s = 6 \tan \theta$. Note that $r^2 + s^2 = (r + s)^2 - 2rs.$ $rs = \frac{6}{\tan \theta} \cdot 6 \tan \theta = 36 \implies 2rs = 72$. The area of the trapezoid is $72 = 6 \cdot \frac{r + s}{2} \implies r + s = 24$. $(r + s)^2 - 2rs = 576 - 72 = \boxed{504}$.

 
                    \begin{figure}[H]
                        \centering
                        \includegraphics[width=0.5\textwidth]{Solution/AIME_2025_I_6_Solution_Figure_1.png}
                    \end{figure}
                    

 -alwaysgonnagiveyouup

 





\textbf{Solution 3 (Fastest formula)}

Denote the radius of the inscribed circle as $R$, and the parallel sides as $r$ and $s$.By formula, we get $R = 3 = \frac{1}{2} \cdot \sqrt{rs}$, where $rs = 36$.Also, by formula, $A = 72 = \frac{1}{2} \cdot \sqrt{rs} \cdot (r + s)$, where $r + s = 24$.Therefore, 
 \begin{align*} &r^2 + s^2 = (r + s)^2 - 2rs \\ &= 24^2 - 2 \cdot 36 \\ &= \boxed{504} \end{align*}


 





\textbf{Solution 4(Double-angle Formula)}

Let \(\angle OAB = \alpha\), \(\tan(\alpha)=\frac{6}{s}\). By the double - angle formula for tangent \(\tan(2\alpha)=\frac{2\tan\alpha}{1-\tan^{2}\alpha}=\frac{2\times\frac{6}{s}}{1 - (\frac{6}{s})^{2}}=\frac{\frac{12}{s}}{\frac{s^{2}-36}{s^{2}}}=\frac{12s}{s^{2}-36}\).

 Since \(\angle DAB = 2\alpha\), \(\tan(2\alpha)=\frac{12}{s - r}=\frac{12s}{s(s - r)}=\frac{12s}{s^{2}-sr}\).

 Set \(\frac{12s}{s^{2}-sr}=\frac{12s}{s^{2}-36}\). Since \(s\neq0\), we can cancel out \(12s\) from both sides of the equation, getting \(s^{2}-sr=s^{2}-36\). Subtracting \(s^{2}\) from both sides, we have \(-sr=-36\), so \(sr = 36\).

 Assume \((r + s)^{2}=576\). Using the formula \((r + s)^{2}=r^{2}+2rs + s^{2}\), then \(r^{2}+s^{2}=(r + s)^{2}-2rs\).

 Substitute \(rs = 36\) and \((r + s)^{2}=576\) into the formula: \(r^{2}+s^{2}=576-2\times36=576 - 72=504\).

 So the final answer is \(504\). By Airbus 320-214.

 Formula reference to here: \href{https://en.wikipedia.org/wiki/Tangential\_trapezoid}{https://en.wikipedia.org/wiki/Tangential\_trapezoid}

 


 





\textbf{Solution 5}

The height of the trapezoid is clearly the diameter of the circle or $6$. We let the larger base be $r$, and the smaller one be $s$. We have by the area of a trapezoid: 
 \[\dfrac{r+s}{2} \cdot 6 = 72 \Longrightarrow r+s = 24\]. Now, by Pitot's theorem, and letting the legs of the trapezoid be $x$, we have that $r+s = 2x = 24 \Longrightarrow x = 12$. Now, by pythag we have that $\dfrac{r-s}{2} = 6\sqrt3 \Longrightarrow r-s = 12\sqrt3$. Now, by systems of equations, we find that $r = 12+6\sqrt3$, and $s = 12-6\sqrt3$. Now, we have that $r^2 + s^2 = (12+6\sqrt3)^2+(12-6\sqrt3)^2 = \boxed{504}$

 -jb2015007

 





\textbf{Solution 6 (Brahmagupta's)}


                    \begin{figure}[H]
                        \centering
                        \includegraphics[width=0.5\textwidth]{Solution/AIME_2025_I_6_Solution_Figure_2.png}
                    \end{figure}
                     

 The area of the trapezoid is $\frac{r + s}{2} \cdot 6 = 72 \implies r + s = 24.$ Note that external tangents of a circle are equal. Let $a$ and $s - a$ be the two external tangents making up the bottom base and $b$ and $r - b$ be the two external tangents making up the top base. Since the trapezoid is isosceles, $a + b = r + s - (a + b) \implies a + b = 12.$ So the legs of the trapezoid have length $12.$ From here, we can apply Brahmagupta's Formula to the trapezoid since isosceles trapezoids are cyclic: $\sqrt{\frac{r + s}{2} \cdot \frac{r + s}{2} \cdot \frac{24 - r + s}{2} \cdot \frac{24 + r - s}{2}} = 72.$ Substituting $r + s = 24$ and solving, we get $(r - s)^2 = 432 \implies s - r = 12\sqrt{3}.$ Since we know $s + r = 24,$ we can just solve this system to get $s = 12 + 6\sqrt{3}$ and $r = 12 - 6\sqrt{3}.$ The answer is $(12 + 6\sqrt{3})^2 + (12 - 6\sqrt{3})^2 = \boxed{504}.$

 





\textbf{Solution 7}

Through Pitot's Theorem,
 \[r+s = 2t\].

 Through tangential quadrilateral area,
 \[3(\frac{r+s+2t}{2}) = 72 \implies 4t = 48 \implies t = 12\].

 Through trapezoidal area,

 
 \[\frac{r+s}{2} \cdot h = 72 \implies h = 6\].

 Through the Pythagorean Theorem or 30-60-90 triangles,

 
 \[m^2 + 36 = 144 \implies m = 6\sqrt{3} \ \text{where} \ 2m + r = s\].

 Through Substitution,

 
 \[r+s = 24, 2r+12\sqrt{3} = 24, r= 12-6\sqrt{3}\].

 Through Backward Substitution,

 
 \[s = 12+6\sqrt{3}\].

 The identity $(a+b)^2 + (a-b)^2 = 2a^2 +2b^2$ forces $r^2 +s^2$ as 

 
 \[2 \cdot 144 + 2 \cdot 108\]
 \[=288 + 216\]
 \[= \boxed{504}\].

 




	\item (\textit{AIME I 2025}) The twelve letters $A$,$B$,$C$,$D$,$E$,$F$,$G$,$H$,$I$,$J$,$K$, and $L$ are randomly grouped into six pairs of letters. The two letters in each pair are placed next to each other in alphabetical order to form six two-letter words, and then those six words are listed alphabetically. For example, a possible result is $AB$, $CJ$, $DG$, $EK$, $FL$, $HI$. The probability that the last word listed contains $G$ is $\frac mn$, where $m$ and $n$ are relatively prime positive integers. Find $m+n$.


 
 \textbf{Answer: }821



\textbf{Solution 1}

Note that order does not matter here. This is because any permutation of the $6$ pairs will automatically get ordered in alphabetical order. The same is true for within each of the pairs. In other words, $AB$ $CH$ $DI$ $EJ$ $FK$ $GL$ should be counted equally as $HC$ $AB$ $DI$ $EJ$ $FK$ $GL$.

 We construct two cases: $G$ is the first letter of the last word and $G$ is the second letter of the last word.

 Our first case is when $G$ is the first letter of the last word. Then the second letter of the last word must be one of $H, I, J, K, L$. Call that set of $5$ letters $\Omega$. There are $5$ ways to choose the second letter from $\Omega$. The other $4$ letters of $\Omega$ must be used in the other $5$ words.

 For the other 5 words, each of their first letters must be before $G$ in the alphabet. Otherwise, the word with $G$ will not be the last. There are $6$ letters before $G$: $A,B,C,D,E,F$. Call that set of $6$ letters $\Sigma$. Exactly one of the words must have two letters from $\Sigma$. The other 4 will have their first letter from $\Sigma$ and the second letter from $\Omega$. There are $4!$ ways to determine the possible pairings of letters from $\Sigma$ and $\Omega$, respectively.

 Therefore, this case has $5 \cdot {6\choose{2}} \cdot 4! = 5 \cdot 15 \cdot 24 = 1800$ orderings.

 The second case is when $G$ is the second letter of the last word. You can see that the first letter of that word must be $F$. Otherwise, that word cannot be the last word. The other $5$ words must start with $A$, $B$, $C$, $D$, and $E$. The second letter of each of those words will come from $\Omega$. There will be $5!$ ways to distribute the elements of $\Omega$ to one of $A, B, C, D, E$. There are therefore $5! = 120$ orderings in the case.

 In total, there are $1800+120 = 1920$ orderings. However, we want the probability. The number of ways to put the $12$ letters into pairs is $11 \cdot 9 \cdot 7 \cdot 5 \cdot 3 \cdot 1$. This is true because we can say this: Start with $A$. It has $11$ options for who it will partner with. There are now $10$ letters left. We must pick the letter that is the least in alphabetical order. It has $9$ options for who it will partner with. There are now $8$ letters left. Continue until there are only $2$ letters left, and there is only $1$ option for that last word. Therefore, there will be  $11 \cdot 9 \cdot 7 \cdot 5 \cdot 3 \cdot 1$ options.

 The probability is therefore $\frac{1920}{11 \cdot 9 \cdot 7 \cdot 5 \cdot 3 \cdot 1} = \frac{128}{693}$. The requested answer is $128 + 693 = \boxed{821}$.

 Minor latex edits by \href{https://artofproblemsolving.com/wiki/index.php?title=User:T3chn0b14d3&action=edit&redlink=1}{T3CHN0B14D3}

 





\textbf{Solution 2: Same but quicker}

Splitting up into $2$ cases: $G$ is the first letter or the second letter of the last word.

 Case $1:$ $G$ in first letter

 Notice that $A$ must take the first letter of first word, one of the letters $B$ - $F$ needs to be the second letter of a word and the rest being the first letter of a word. The combinations will be $1 + 2 + 3 + 4 + 5 = 15.$ After the first $7$ letters has been decided then the last $5$ will just fill by $5!.$ This case will have $15 \cdot 5!$ outcomes.

 
Case $2:$ $G$ in last letter

 Notice that $A$ - $G$ has been arranged by $A? B? C? D? E? FG,$ where the $?$ is undecided. We have another $5!$ to fill out the possible outcomes. 

 In total, there are $16 \cdot 5!.$ The total case will be $11 \cdot 9 \cdot 7 \cdot 5 \cdot 3 \cdot 1$ (Consider A must be in the first letter of first word, then you have $11$ choices, then you must take the next letter in alphabetical order as mandatory, then you have a free choice of $9$ and so on). 

 Answer:
 \[= \frac{16 \cdot 5 \cdot 4 \cdot 3 \cdot 2 \cdot 1}{ 11 \cdot 9 \cdot 7 \cdot 5 \cdot 3 \cdot 1}\]
 \[= \frac{16 \cdot 4 \cdot 2}{11 \cdot 9 \cdot 7}\]
 \[= \frac{128}{ 693}\]Therefore it gives us the answer of ${128 + 693 = \boxed{821}.}$

 \textbf{Alternative Step in Solutions 1 and 2: Total Arrangements} The total number of possible arrangements can also be calculated using combinations. We wish to find the total number of ways to form $6$ pairs of letters from the given $12$. There are $12!$ ways to do so. However, the $2$ letters in each pair can only be arranged in $1$ way, not $2!$ ways, because they must be in alphabetical order. Thus, we must divide by $2!$ for each of the $6$ pairs, giving $\frac{12!}{2!2!2!2!2!2!}$ or $\frac{12!}{2^6}$ arrangements. 

 
We must also take into account the number of ways to arrange the $6$ pairs relative to each other. In the problem, only $1$ possible arrangement can arise from $6!$ different cases because the pairs of letters must be in alphabetical order. Therefore, we must divide by $6!$ to give the total number of possible arrangements as $\frac{12!}{2^6\cdot6!}=\frac{2^{10}\cdot3^5\cdot5^2\cdot7\cdot11}{2^6\cdot2^4\cdot3^2\cdot5}=3^3\cdot5\cdot7\cdot11=10,395$.

 \textbf{Alternative view 2} Similar to what he said above, we can do $\frac{\binom{12}{2} * \binom{10}{2} * \binom{8}{2} * \binom{6}{2} * \binom{4}{2} * \binom{2}{2}}{6!}$

 This makes 6 pairs and we divide by 6! because there is only one configuration which will be in alphabetical order

 





\textbf{Solution 3 (Rigorous)}

Propose the set $\mathcal{L}$ that determines all plausible cases. Within the total, we have the opportunity to arrange the 12 letters in $12!$ ways and split it into groups of 2. Within each group of 2, we have 2 permutations, only one of which is chosen, so that is $\frac{12!}{2^6}$. Now, within the total collective group, out of $6!$ arrangements, only 1 is deemed viable. Therefore, the total is $\frac{12!}{2^6 \cdot 6!} = 10395$ cases.

 Now, propose $\mathbf{\Omega} \in \mathcal{L}$ that persuades $G$ as the first letter of the last word. The remaining letter in $G$ must be chosen from $\{H, I, J, K, L \}$ in 5 ways. 

 Without loss of generality, assume $GH$ proposes the last word. Consider the last five words chosen from the set $\{A,B,C,D,E,F \} \cup \{I,J,K,L \}$. We purpose a similar solution to solution 1. It is evident because six balls $A,B,C,D,E$ are to be dropped in four boxes with two remainder, those two balls must pair with one another as a result by the Pigeonhole Principle. We have $\binom{6}{4}$ ways to drop the four letters within the boxes. However, we have the sole exception of the number of ways to order said boxes. Yes, we can choose four elements from the set to become ordered, but we have choices on what to order them to. Obviously, every one of the elements has $4$ options, then the remainder has 3, then 2, then 1, yielding $4! = 24$ cases. The total number of partitions here is therefore $24 \cdot 5 \cdot 15 = 1800$.

 Propose $\mathbf{\Phi} \in \mathcal{L}$ that assumes $G$ as the last letter of the last word. The property of degeneracy tells us that if any letter of the previous set of $\mathbf{\Omega}$ is chosen to represent a letter in set $\mathbf{\Phi}$ then $\mathbf{\Phi} \cap \mathbf{\Omega} = \mathbf{\Omega}$, which is a degeneracy condition, and therefore we require $\mathbf{\Phi}$ to be distinctive. So, $G \in \{H,I,J,K,L \}^C \implies G \in \{A,B,C,D,E,F \}$. Propose $\mathbf{\Gamma} = \mathcal{L}^C$. Then, for $G \in \{A,B,C,D,E \}$, the property of degeneracy tells us that $\mathbf{\Gamma} \cap \mathbf{\Phi} = \emptyset$. Therefore, $G \in \{F \}$ so the term $FG$ is enumerated as the last term.  

 Now consider the last five words to be formed of letters, proposed from the set $\{A,B,C,D,E \} \cup \{H,I,J,K,L \}$. We put each of five balls in five boxes and through Pigeonhole Principle, that leaves 0 balls remaining. Within such groups, we have $5! = 120$ distinctive orderings, conjoining a total 120 cases.

 We request the probability deterministic of
 \[\frac{\mathbf{\Phi} \cup \mathbf{\Omega}}{\mathcal{L}}\]which is exactly
 \[\frac{1920}{10395}\]which furthermore simplifies to 
 \[\frac{128}{693} \implies \boxed{821}\].

 The rigorous proof of the count has been complete.

 




	\item (\textit{AIME I 2025}) Let $k$ be a real number such that the system

 \begin{align*} &|25 + 20i - z| = 5 \\ &|z - 4 - k| = |z - 3i - k| \end{align*}
has exactly one complex solution $z$. The sum of all possible values of $k$ can be written as $\frac{m}{n}$, where $m$ and $n$ are relatively prime positive integers. Find $m + n$. Here $i = \sqrt{-1}$.


 
 \textbf{Answer: }077



\textbf{Solution 1}


                    \begin{figure}[H]
                        \centering
                        \includegraphics[width=0.5\textwidth]{Solution/AIME_2025_I_8_Solution_Figure_1.png}
                    \end{figure}
                    The complex number $z$ must satisfy the following conditions on the complex plane: 

 $1.$ The magnitude between $z$ and $(25,20)$ is $5.$ This can be represented by drawing a circle with center $(25,20)$ and radius $5.$

 $2.$ It is equidistant from the points $(4+k,0)$ and $(k,3).$ Hence it must lie on the perpendicular bisector of the line connecting these points. 

 
For $z$ to have one solution, the perpendicular bisector of the segment connecting the two points must be tangent to the circle. This bisector must pass the midpoint, $(2+k,\frac{3}{2}),$ and have slope $\frac{4}{3}.$ The segment connecting the point of tangency to the center of the circle has slope $\frac{-3}{4},$ meaning the points of tangency can be $(29,17)$ or $(21,23).$ Solving the equation for the slope of the perpendicular bisector gives 
 \[\frac{\frac{3}{2}-23}{k+2-21}=\frac{4}{3}\] or 
 \[\frac{\frac{3}{2}-17}{k+2-29}=\frac{4}{3},\] giving $k=\frac{23}{8}$ or $\frac{123}{8}$, having a sum of $\frac{73}{4} \Longrightarrow \boxed{077}.$

 





\textbf{Solution 2 (Systematic + Algebra)}

We first look at each equation, and we convert each to algebra (note that the absolute value sign of $|$ means the magnitude). Let's convert z to $A + Bi$.

 
Note that the first equation becomes:$(25 - A)^2 + (20 - B)^2 = 25$

 Note that this is the equation of a circle centered at $(25, 20)$ with radius $5$.

 
And the second equation becomes:$(A-4-k)^2 + B^2 = (A - k)^2 + (B-3)^2$

 You can see that the many similar terms that cancel out, simplfying, you get:

 $-8(A - k) + 16 + 6B = 9$

 Now we must isolate B

 
 \[B= \frac{4}{3}(A-k) - \frac{7}{6}\]

 
 \[B = \frac{4}{3}A - \frac{4}{3}k - \frac{7}{6}\]

 This equation can be seen as a line with a $\frac{4}{3}$ slope, and a y-intercept of $- \frac{4}{3}k - \frac{7}{6}$. 

 Note that the question only wants one solution, so we want two tangent lines, one above the circle, and one below the circle. You can see Solution 2 for the picture.

 Because the slope is $\frac{4}{3}$, the circle must have a slope coming out of center of its reciprocal, $-\frac{3}{4}$. So the points on the circle where this line with a $\frac{4}{3}$ must intersect must be $(21, 23)$ and $(29, 17)$. We can easily use point-slope form to find the equations of these lines. 
 \[y - 23 = \frac{4}{3}(x - 21)\]

 and 

 
 \[y - 17 = \frac{4}{3}(x - 29)\]

 Now we must match the y-intercepts to the equations with $k$ in it.Solving the equations: 

 
 \[\frac{4}{3}(-21) + 23 = - \frac{4}{3}k - \frac{7}{6}\]

 
 \[\frac{4}{3}(-29) + 17 = - \frac{4}{3}k - \frac{7}{6}\]

 we get that $k = \frac{23}{8}$ and $k = \frac{123}{8}$Adding them up and simplifying, we get a sum of $\frac{73}{4} \Longrightarrow \boxed{077}.$

 





\textbf{Solution 3 (A Little Geometry)}


                    \begin{figure}[H]
                        \centering
                        \includegraphics[width=0.5\textwidth]{Solution/AIME_2025_I_8_Solution_Figure_2.png}
                    \end{figure}
                    
                    \begin{figure}[H]
                        \centering
                        \includegraphics[width=0.5\textwidth]{Solution/AIME_2025_I_8_Solution_Figure_3.png}
                    \end{figure}
                    

 To solve the problem, we first locate the point $Z$. According to the conditions, we can know that:

 $Z$ is on the perpendicular bisector of $(k,3)$ and $(k+4,0)$

 The distance from $Z$ to circle $O(25,20)$ is ${5}$.

 Therefore, $Z$ must be the 2 points of tangent of a line with the slope of $\frac{4}{3}$ with circle O (center at $(25,20)$, radius of ${5}$), corresponding to the 2 values: ${K_1}$ and ${K_2}$.

 Since the question only asks for the sum of ${K_1}$ and ${K_2}$, we would not need to calculate them separately. Let the middle point of ${K_1}$ and ${K_2}$ be ${K}$ ==> ${K_1}$ + ${K_2}$ = ${2K}$.

 Looking at the simplified figure below. We may calculate using similarity of the triangles:

 $\frac{SV}{SO}=\frac{WV}{OT}$

 $SV = \frac{2.5\cdot25}{20}=\frac{25}{8}$

 $KS = KV-SV = 4-\frac{25}{8} = \frac{7}{8}$

 $KT = KS+ST = \frac{7}{8}+15 = \frac{127}{8}$

 $K_1+K_2 = 2K = (25-\frac{127}{8})\cdot2= \frac{73}{8}\cdot2 = \frac{73}{4}$

 We conclude that $m+n=\boxed{077}$.

 





\textbf{Solution 4 (Distance Formula, quick)}

Let $z = a + bi.$ Then the system becomes $(25 - a)^2 + (20 - b)^2 = 25$ and $(a - 4 - k)^2 + b^2 = (a - k)^2 + (b - 3)^2$ $\implies b = \frac{4}{3}a - (\frac{4}{3}k + \frac{7}{6})$. For this system to have one solution for $z,$ the line $b = \frac{4}{3}a - (\frac{4}{3}k + \frac{7}{6})$ must be tangent to the circle $(25 - a)^2 + (20 - b)^2 = 25$. For a line to be tangent to a circle, the distance between the line and center must be equal to the length of the radius, which is $5$. The center of the circle is $(a,b) \rightarrow (25,20),$ and we can express the line as $8a - 6b - (8k + 7) = 0.$ So we have: 
 \[d = \frac{|Ax_0 + By_0 + C|}{\sqrt{A^2 + B^2}} = \frac{|8(25) - 6(20) -(8k + 7)|}{\sqrt{8^2 + 6^2}} =  \frac{|73 - 8k|}{10} = 5 \implies |73 - 8k| = 50\] 
 \[k = \frac{23}{8}, \frac{123}{8}.\] $\frac{m}{n}$ is $\frac{146}{8} = \frac{73}{4}.$ The requested sum is $73 + 4 = \boxed{77}.$

 





\textbf{Solution 5 (Educated Guess)}

Letting $z=a+bi$, the given equations become $(a-25)^2+(b-20)^2 = 25$ and $(a-4-k)^2+b^2 = (a-k)^2+(b-3)^2$. Expanding the second equation and canceling like terms yields $16-8a+8k=-6b+9$. Rearrange and solve for $k$ to get $k=\frac{8a-6b-7}{8}$.

 The equation $(a-25)^2+(b-20)^2 = 25$ is a circle with center $(25, 20)$ and radius $5$. As explained in Solution 1, $|z - 4 - k| = |z - 3i - k|$ represents the perpendicular bisector of the line segment formed by connecting points $(4+k, 0)$ and $(k, 3)$. The $2$ tangent points hit on the circle by the perpendicular bisector with be diametrically opposite.

 Experiment with plugging in diametrically opposite pairs of points such as $(a, b)={(20, 15), (20, 25)}$ and $(a, b)={(29, 17), (21, 23)}$, which can be found using $3-4-5$ right triangles, into the previously found formula $k=\frac{8a-6b-7}{8}$. Observe that when the resulting $k$-values from each pair are added together, the value $\frac{146}{8}=\frac{73}{4}$ is obtained every time. If one assumes that this holds true for each pair of points, then it also holds true for the $2$ desired points, and the answer is $73+4=\boxed{077}$.

 




	\item (\textit{AIME I 2025}) The parabola with equation $y = x^2 - 4$ is rotated $60^{\circ}$ counterclockwise around the origin. The unique point in the fourth quadrant where the original parabola and its image intersect has $y$-coordinate $\frac{a - \sqrt{b}}{c}$, where $a$, $b$, and $c$ are positive integers, and $a$ and $c$ are relatively prime. Find $a + b + c$.


 
 \textbf{Answer: }062



\textbf{Solution 1}

We need to relate the rotation to something simpler because substituting the equation of the rotated parabola into the original will give us a quartic. 
                    \begin{figure}[H]
                        \centering
                        \includegraphics[width=0.5\textwidth]{Solution/AIME_2025_I_9_Solution_Figure_1.png}
                    \end{figure}
                    

 Notice that the vertices of the original parabola and its image are symmetrical about the angle bisector of the 60 degree rotation (shown in green).

 The equation of this line is $y = -\tan60^\circ \cdot y = -\sqrt{3}x.$ This means that the point lying on this line and the original parabola must be the intersection of the new and original images since it won't change position with the rotation. 

 
We substitute $y = x^2 - 4:$ 
 \[x^2 - 4 = -\sqrt{3}x\] 
 \[x = \frac{-\sqrt{3} + \sqrt{19}}{2}.\] 

 Then $y = (\frac{-\sqrt{3} + \sqrt{19}}{2})^2 - 4 = \frac{3 - \sqrt{57}}{2}.$ $57 + 3 + 2 = \boxed{62}.$

 \textbf{Note} One can find the remaining two intersection points using the graph $y = \frac{1}{\sqrt{3}}x - \frac{2}{3}$ and $y=x^2 - 4$.

 \href{https://www.desmos.com/calculator/rlqc4ikrc6}{https://www.desmos.com/calculator/rlqc4ikrc6}

 





\textbf{Solution 2 (Similar to Solution 1)}

Note that this question is equivalent to finding a point $B$ in the fourth quadrant, such that when a point $A$ on the graph of $y = x^2 - 4$ is rotated $60^\circ$ counterclockwise around the origin, it lands on $B$, which is also on the graph.

 The first thing to note is that point $A$ and $B$ must be equidistant to the origin. If we express the coordinates of $A$ as $(a, b)$, and the coordinates of $B$ as $(x, y)$, we have:
 \[\|OA\| = \|OB\|\]which means that:
 \[\sqrt{a^2 + b^2} = \sqrt{x^2 + y^2}\]Since $b = a^2 - 4$ and $y = x^2 - 4$, we have $a^2 = b + 4$ and $x^2 = y + 4$, substituting this into the previous equation and squaring both sides yields:
 \[2a^2 + 4 = 2x^2 + 4\]Meaning that $a^2 = x^2$, since $A$ and $B$ clearly cannot coincide, we must have $a = -x$, since $y = x^2 - 4$ is an even function, this means that point $A$ and $B$ are just reflections of each other over the y axis. The angle between $\overline{OA}$ and $\overline{OB}$ is $60^\circ$ and $A$ and $B$ is symmetrical, the y axis should bisect the angle \angle AOB, i.e., the angle between $\overline{OB}$ and the y axis is:
 \[\frac{60^\circ}{2} = 30^\circ\]  Therefore the point $B$ must lie on the line 
 \[y = -\sqrt{3}x\]

 We have:
 \[\begin{cases}y = x^2 - 4 \\ y = -\sqrt{3}x \end{cases}\]
 \[x^2 - 4 = -\sqrt{3}x\]Using the quadratic formula and keeping in mind that the x value is positive (since $B$ is in the fourth quadrant) yields $x = \frac{\sqrt{19} - \sqrt{3}}{2}$.

 Substituting into 
 \[y = -\sqrt{3}x\] We get 
 \[y=\frac{3-\sqrt{57}}{2}\implies \boxed{062}.\]

 The last part of this solution is essentially Solution 1.

 The last part of this solution is essentially Solution 1.

 





\textbf{Solution 3 (clarification needed): Rotation}

Rotate the point $(x,y)$ $60^\circ$ counterclockwise about the origin. To find the resulting point $(x',y')$, use polar coordinates. Let $x=r \cos\theta$ and $y=r \sin\theta$, so $(x',y')=(r \cos({\theta+60^\circ}), r \sin({\theta+60^\circ}))$. Expanding using the angle addition identities and simplifying gives $(x',y')=\left(\frac{r\cos\theta-r\sqrt{3}\sin\theta}{2}, \frac{r\sin\theta+r\sqrt{3}\cos\theta}{2}\right)$. Substituting back in $x=r \cos\theta$ and $y=r \sin\theta$ yields $(x',y')=\left(\frac{x-y\sqrt{3}}{2}, \frac{x\sqrt{3}+y}{2}\right)$.

 
To rewrite $(x',y')$ completely in terms of $x$, substitute $y=x^2-4$ to get $(x',y')=\left(\frac{x-(x^2-4)\sqrt{3}}{2}, \frac{x\sqrt{3}+(x^2-4)}{2}\right)$.) 

 
We now have the coordinates for each point on the rotated parabola: $(x',y')=\left(\frac{x-(x^2-4)\sqrt{3}}{2}, \frac{x\sqrt{3}+(x^2-4)}{2}\right)$. Set the $y$-coordinates for both the original and rotated parabola equal to each other: $\frac{x\sqrt{3}+(x^2-4)}{2}=x^2-4$. Rearranging gives the quadratic $x^2-\sqrt{3}x-4=0$, and applying the Quadratic Formula yields the solutions $x=\frac{\sqrt{3}\pm\sqrt{19}}{2}$. 

 
Taking the solution $x=\frac{\sqrt{3}-\sqrt{19}}{2}$ (IDK why though, the $x$-value is negative when it should be positive, clarification needed) (due to the way we set up the equation, we are finding the $x$ value of the point that maps to the intersection point after rotation. $-x$ is the $x$ value for the intersection point ~Towerrocks)  and substituting it into $y=x^2-4$ gives $y=\left(\frac{\sqrt{3}-\sqrt{19}}{2}\right)^2-4=\frac{22-2\sqrt{57}}{4}-4=\frac{6-2\sqrt{57}}{4}=\frac{3-\sqrt{57}}{2}$. Therefore, our answer is $a+b+c=3+57+2=\boxed{062}$.

 


 Note: It may be possible to solve the problem by substituting the coordinates $(x',y')=\left(\frac{x-y\sqrt{3}}{2}, \frac{x\sqrt{3}+y}{2}\right)$ into the equation $y=x^2-4$ to obtain the systems of equations $\frac{x\sqrt{3}+y}{2}=\left(\frac{x-y\sqrt{3}}{2}\right)^2-4, y=x^2-4$.

 \href{https://mathproblemsolvingskills.wordpress.com/}{MathProblemSolvingSkills.com}

 




	\item (\textit{AIME I 2025}) The $27$ cells of a $3 \times 9$ grid are filled in using the numbers $1$ through $9$ so that each row contains $9$ different numbers, and each of the three $3 \times 3$ blocks heavily outlined in the example below contains $9$ different numbers, as in the first three rows of a Sudoku puzzle. 


 
                    \begin{figure}[H]
                        \centering
                        \includegraphics[width=0.5\textwidth]{Problem/AIME_2025_I_10_Problem_Figure_1.png}
                    \end{figure}
                    


 The number of different ways to fill such a grid can be written as $p^a \cdot q^b \cdot r^c \cdot s^d$ where $p$, $q$, $r$, and $s$ are distinct prime numbers and $a$, $b$, $c$, $d$ are positive integers. Find $p \cdot a + q \cdot b + r \cdot c + s \cdot d$.


 
 \textbf{Answer: }081



\textbf{Solution 1}

We will fill out the grid row by row. Note that there are $9!=2^7\cdot 3^4\cdot 5\cdot 7$ ways to fill out the first row. For the second row, we will consider a little casework. WLOG let the first row be $123 | 456 | 789$ (bars indicate between the $3\times 3$ squares).

 <u>Case 1:</u> Every number from the first square goes to the second square. Then there is $1$ way along with $6^3$ orderings giving us $6^3$ cases.

 <u>Case 2:</u> Two of the first numbers go in the second square and the last number goes in the third. This results in $3\cdot3\cdot3$ ways to put the numbers along with $6^3$ orderings. Thus there are $27\cdot 6^3$ cases.

 Note that we didn't consider $2$ cases but these are symmetric so the number of ways to fill the second row is $56\cdot 6^3$.

 There is $1$ way to place the last numbers in with $6^3$ orderings. Thus there are $2^7\cdot 3^4\cdot 5\cdot 7 \cdot 56\cdot 6^3 \cdot 6^3 = 2^{16}\cdot 3^{10}\cdot 5^1\cdot 7^2$ ways to fill out our mini Sudoku puzzle. Thus the answer is $2\cdot 16+3\cdot 10 + 5\cdot 1+7\cdot 2 = \boxed{081}.$

 





\textbf{Solution 2}

We will call the sudoku blocks 1-3 from left to right, and rows 1-3 from up to down. Lets start from the entire first block - we have $9!$ options of choosing it as you just arrange the numbers.

 Starting from the second block, we need to do some casework. We need to discuss the different possibilities that the first row of the second block contains. 

 Case 1: The first row of the second block has the exact 3 numbers from the second row of first block. 

 We have $P^3_3$ choices as you just arrange the numbers for the first row. 

 Then, considering the second row of the second block, you must take the exact 3 numbers from the third row of first block, and we could arrange then in different order. This will gives another $P^3_3$ choices. 

 The third row of the second block must contain the 3 numbers of the first row in first block. We will then have another $P^3_3$ of arranging the numbers. 

 Summing up case 1, we have $P^3_3 \cdot P^3_3 \cdot P^3_3 = 6^3 = 2^3 \cdot 3^3$.

 Case 2: The first row of second block contains 2 numbers from the second row of first block and 1 number from the third row of first block.

 To satisfy my case 2 statement, we will need $\binom{3}{1}\cdot\binom{3}{2}\cdot P^3_3$.

 To look at the second row of the second block, we must first choose the remaining 2 numbers that haven't got picked from the third row of first block, or else the sudoku won't work. This means that we can only pick 1 remaining number from the numbers in first row of first block - because the 6th number is the number in the second row and sudoku doesn't allow repeated numbers in a row. Therefore, the second row has the choice of $\binom{2}{2}\binom{3}{1}P^3_3$. 

 The third row of the second block gets the remaining 3 numbers that is left from the second block, which is $P^3_3$. 

 Summing up case 2, we have $\binom{3}{1}\binom{3}{2}P^3_3\binom{2}{2}\binom{3}{1}P^3_3P^3_3= 3^3\cdot2^3\cdot3^3 = 2^3\cdot3^6$.

 Notice that case 3 (3 numbers repeating from the third row of first block) and case 4 (2 numbers from block 1 row 3 and 1 number from block 1 row 2) are symmetric to case 1 and 2, respectively. Therefore the second block results with $(1 + 3^3 + 1 + 3^3) \cdot 2^3 \cdot 3^3 = 56 \cdot 2^3 \cdot 3^3 = 7 \cdot 2^6 \cdot 3^3$

 Third block: easy

 Notice that the 6 numbers of each row has been filled already, the choices of the rest 3 numbers has been done. You only need to arrange them in different orders, which gives $P^3_3P^3_3P^3_3= 2^3 * 3^3$.

 In total, our answer is:

 $9! \cdot 7 \cdot 2^6 \cdot 3^3 \cdot 2^3 \cdot 3^3$$=2 \cdot 3 \cdot 2^2 \cdot 5 \cdot 2 \cdot 3 \cdot 7 \cdot 2^3 \cdot 3^2 \cdot 7 \cdot 2^9 \cdot 3^6$$=2^{16} \cdot 3^{10} \cdot 5 \cdot 7^2$

 Our required answer is therefore $2 \cdot 16 + 3 \cdot 10 + 5 \cdot 1 + 7 \cdot 2 = 32 + 30 + 5 + 14 = \boxed{081}$

 Minor latex edits by \href{https://artofproblemsolving.com/wiki/index.php?title=User:T3chn0b14d3&action=edit&redlink=1}{T3CHN0B14D3}

 





\textbf{Solution 3 (Can be very quick but explained well)}

First, assume the first row is ordered $1, 2, 3, 4, 5, 6, 7, 8, 9$. 

 This contributes $9!$ cases. 

 Next, assume the first box has the remaining numbers, $4, 5, 6, 7, 8, 9$ ordered in some way. This contributes $6! = 720$ cases for the first box.

 Third, consider each of the remaining $1$ x $3$ boxes as separate, and they each contribute $3!$ cases. 

 However, we are undercounting. 

 The second box may have to be multiplied by 3 if the first box did not define the numbers that go in each $1$ x $3$ box. The third box will always have its rows defined. If the first box is placed such that $4, 5, 6$ are in the same $1$ x $3$ box, and $7, 8, 9$ are in the same $1$ x $3$ box, then the numbers in the second box will be defined and we do not need to multiply by 3. This happens for $2 \cdot 3! \cdot 3! = 72$ cases, or $72 / 720 = 1 / 10$ of the time. Therefore, $9 / 10$ of the time, we need to multiply by another 3. Adding this up, to account for undercounting, we must multiply the total by $1 / 10 + 3(9 / 10) = 28 / 10 = 14 / 5$.

 Our final total is $9! \cdot 6! \cdot (3!)^4 \cdot 14 / 5 = (2^7 \cdot 3^4 \cdot 5^1 \cdot 7^1) (2^4 \cdot 3^2 \cdot 5^1) (2^4 \cdot 3^4) \cdot 2 \cdot 7 / 5 = 2^{16} \cdot 3^{10} \cdot 5^1 \cdot 7^2$. 

 The desired solution is $2\cdot16 + 3\cdot10 + 5\cdot1 + 7\cdot2 = \boxed{081}$

 Minor latex edits by Wen\_Liang

 





\textbf{Solution 4}

Assume the first row of the grid is $123|456|789.$ Note that the numbers in each $3$x$1$ block can be ordered in $3!$ ways amongst themselves and they will still stay in the same row and $3$x$3$ block. There are $6$ unfilled $3$x$1$ blocks, so we'll multiply by $(3!)^6$ at the end. We’ll also have to multiply by $9!$ because we fixed the first row of the grid and there are $9!$ ways to permute the numbers in it.

 Let's do casework on where $789$ can be in the second row of the grid. 

 \textbf{Case 1:} All three numbers $789$ are in the second row of the same $3$x$3$ block 

 $789$ can be the second row of one of two $3$x$3$ blocks (it can't be the second row of the third $3$x$3$ block since the top row of that one already has $789$ in it). WLOG say $789$ is the second row of the first $3$x$3$ block. Then $123$ must be the second row of the second $3$x$3$ block and $456$ must be the second row of the third $3$x$3$ block. Since the second row is completely filled out we are done. So this case gives us 2 ways. 

 
\textbf{Case 2:} Two of the numbers of $789$ are in the second row of one $3$x$3$ block, the other left out number is in the second row of the other $3$x$3$ block

 We can choose the "left out" number in $\binom{3}{1} = 3$ ways, and it can go in the second row of one of two blocks (remember, it can't be in the third $3$x$3$ block since the top row of that one already has $789$ in it). So we have $3\cdot 2 = 6$ ways. WLOG say the left out number is $9,$ and we put $9$ in the second row of the first $3$x$3$ block, so $7$ and $8$ are in the second row of the second $3$x$3$ block. Now, consider the numbers $456.$ We can choose two of these numbers to join $9$ in the second row of the first $3$x$3$ block in $\binom{3}{2} = 3$ ways. WLOG say we chose $4$ and $5$. Then $6$ must be in the second row of the third $3$x$3$ block. Now we can choose one number out of $1,2,3$ to join $7$ and $8$ in the second row of the second block in $\binom{3}{1} = 3$ ways. WLOG say we chose $1.$ Then $2$ and $3$ must be in the second row of the third $3$x$3$ block with $6.$ Now the second row is completely filled out, so we are done. We have $6 \cdot 3 \cdot 3 = 54$ ways for this case. 

 The total number of ways is $(54 + 2) \cdot (3!)^6 \cdot 9! = 2^{16} \cdot 3^{10} \cdot 5^1 \cdot 7^2,$ so the answer is $32 + 30 + 5 + 14 = \boxed{81}.$

 





\textbf{Solution 5 (More in depth sol 1)}

There are $9!$ ways to order the first row. If we choose the second row, the third row is uniquely determined.

 Therefore, do casework on second row.

 Case 1: All elements in the top $1 \times 3$ grid of the first grid is given to the second. Then, the top $1 \times 3$ grid of the second must be given to the third and the same for the third. That contributes one such ordering. However, within the row we have $3! = 6$ permutes each, accomplishing $6^3$ different choices.

 Case 2: Two of three in the top $1 \times 3$ grid of the first is given to the second and the last for the third. There are 3 ways to choose which number goes to grid 3. Now, from the second grid, we must give one number to the first grid and two to the third, as if we did it vice versa, grid 3 breaks the sudoku. Therefore, grid three follows the same way, and that is 27 cases total. Each row has 6 permutations, so there are $27 \cdot 6^3$ ways.

 Case 3: One of three in the top $1 \times 3$ grid of the first is given to the second and the rest for the third. This is symmetric to case 2, so also $27 \cdot 6^3$ ways.

 Case 4: All elements in the top $1 \times 3$ grid of the first is given to the third. Symmetric to case 1, so $6^3$ ways.

 That is a total $56 \cdot 6^3$ ways, and so multiplying with $9!$ and taking the prime factorization through a mix of Legendre and Bashing, we have $2^{16} \cdot 3^{10} \cdot 5^1 \cdot 7^2 \implies \boxed{81}$.

 





\textbf{Solution 6(Similar to above solutions but with a little more explanation)}

Note that the first row of the grid has to be some permutation of the numbers $(1, 2, \cdots, 9)$ since that's how a Sudoku works. WLOG, assume the first row is $1, 2, 3, \cdots, 9$ in that order. At the end, since the first row could be any permutation of this set, we will multiply by $9!$. 

 Now observe that once we fill out the second row of the grid, we can start looking at filling out each of the three individual boxes. Since the first and second rows are already filled out, the third row of each individual box already is determined by the numbers filled out in the first and second rows of that box. Thus, it's just a matter of the permutations of that third row. For example, let's say we take the first box and as the first row is filled out, let's say the first row is $(1, 4, 7)$. Then let's say we fill out the second row as $(2, 6, 9)$. Then the third row of that box must only be the leftover three numbers namely $(3, 5, 8)$ and we have $3!$ ways to arrange these numbers for the last row. Since this happens for all three of the individual boxes, we can multiply by $(3!)(3!)(3!) = (3!)^3$ at the end.

 Thus, we will multiply by $9! \cdot (3!)^3$ at the end. 

 Now, it just remains to find the number of ways to fill out the second row given that the first row is already filled out. Consider the first row of the first individual box. It is $(1, 2, 3)$ from our WLOG assumption. Thus, the second row has to be a $3$-element subset of the remaining numbers namely $(4, 5, 6, 7, 8, 9)$. Now consider the first row of the second individual box which is $(4, 5, 6)$. Similarly the first row of the third individual box is $(7, 8, 9)$. We will denote these $3$-element subsets as groups. Let $(1, 2, 3)$ be group $A$. Similarly, let $(4, 5, 6)$ be group $B$ and $(7, 8, 9)$ be group $C$. 

 For the second row of the first individual box, as we established earlier, it is any $3$-element subset of $(4, 5, 6, 7, 8, 9)$ which is exactly the union of group $B$ and group $C$. We can do some casework.

 Case 1: To form the $3$-element subset for the second row of the first box, all $3$ elements come from group $B$.

 There is exactly $\binom{3}{3}$ ways to select all $3$ elements from group $B$. Then the second row of the second box must come from $(1, 2, 3, 7, 8, 9)$ which is exactly the union of group $A$ and group $C$. However, note that if we use group $A$ to fill out the second row of the second box, then it remains for group $C$ to be used to fill out the second row of the third box; however that will create the same numbers appearing twice in the same row which is a contradiction. Thus, we must only use group $C$ to fill out the second row of the second box and then it is left for group $A$ to finish the job. Thus, there are $\binom{3}{3}$ to use the $3$-elements from group $C$ and same for group $A$. Note that the order of the numbers appearing in each row of the box matters so we multiply by $3!$ to arrange the elements in the first box, and similarly the same for the second and third boxes giving $\binom{3}{3} \cdot \binom{3}{3} \cdot \binom{3}{3} \cdot 3! \cdot 3! \cdot 3! = (3!)^3$.

 Case 2: To form the $3$-element subset for the second row of the first box, we take exactly two elements from group $B$ and one element from group $C$.

 Clearly we have $\binom{3}{2}$ ways for group $B$ and $\binom{3}{1}$ ways for group $C$. WLOG, assume we take $(4, 5)$ from group $B$ and $(7)$ from group $C$ leaving the numbers left to fill out the second row of the second box to be

 $(5, 8, 9, 1, 2, 3)$ 

 Note that to prevent using the same numbers twice for the second row of the third box, we need to use up $(8, 9)$ for the second box and additionally $5$ cannot be used since that will create $5$ appearing twice. So we have to choose exactly $1$ remaining element from $(1, 2, 3)$ giving $\binom{3}{1}$. Again, we multiply by $(3!)^3$ since the order of the numbers appearing in each row of the boxes matters. Thus, we get

 $\binom{3}{2} \cdot \binom{3}{1} \cdot \binom{3}{1} \cdot (3!)^3 = 27 \cdot (3!)^3$.

 Case 3: To form the $3$-element subset for the second row of the first box, we take exactly one element from group $B$ and two elements from group $C$.

 This will yield $\binom{3}{1}$ ways for group $B$ and $\binom{3}{2}$ ways for group $C$. WLOG, assume we take $(4)$ from group $B$ and $(7, 8)$ from group $C$. Thus, the remaining numbers that can be used to fill out the second row of the second box are $(5, 6, 9, 1, 2, 3)$. But again, to prevent double numbers in the third box, we need to already use up the $(9)$ and we cannot use $(5, 6)$ since then $5, 6$ will appear twice in the second box. Thus, we need to select exactly two remaining elements from $(1, 2, 3)$ giving $\binom{3}{2}$. Remember to multiply everything by $(3!)^3$ since order matters.

 Thus we get $\binom{3}{1} \cdot \binom{3}{2} \cdot \binom{3}{2} \cdot (3!)^3 = 27 \cdot (3!)^3$.

 Case 4: To form the $3$-element subset for the second row of the first box, all $3$ elements come from group $C$. 

 This has $\binom{3}{3}$ ways for group $C$. Then the numbers remaining to fill out the second row of the second box are

 $(4, 5, 6, 1, 2, 3)$. Again, by similar logic to Case 1, we can only use group $A$ to fill out the second row of the second box leaving group $B$ to fill out the second row of the third box which is the only way it will work out(otherwise there will be double numbers). Again, remember to multiply by $(3!)^3$. Thus, we get $\binom{3}{3} \cdot \binom{3}{3} \cdot \binom{3}{3} \cdot (3!)^3 = (3!)^3$.

 Finally, adding up all the cases gives

 $(3!)^3 + 27 \cdot (3!)^3 + 27 \cdot (3!)^3 + (3!)^3 = (1 + 27 + 27 + 1)(3!)^3 = 56 \cdot (3!)^3$.

 Then, remember we established to multiply everything by $9! \cdot (3!)^3$ so the expression becomes

 $9! \cdot (3!)^3 \cdot 56 \cdot (3!)^3 = 2^{16} \cdot 3^{10} \cdot 5^{1} \cdot 7^{2} \implies \boxed{081}$

 




	\item (\textit{AIME I 2025}) A piecewise linear function is defined by 
 \[f(x) = \begin{cases} x & \operatorname{if} ~ -1 \leq x < 1 \\ 2 - x & \operatorname{if} ~ 1 \leq x < 3\end{cases}\] and $f(x + 4) = f(x)$ for all real numbers $x$. The graph of $f(x)$ has the sawtooth pattern depicted below.


 
                    \begin{figure}[H]
                        \centering
                        \includegraphics[width=0.5\textwidth]{Problem/AIME_2025_I_11_Problem_Figure_1.png}
                    \end{figure}
                    


 The parabola $x = 34y^{2}$ intersects the graph of $f(x)$ at finitely many points. The sum of the $y$-coordinates of all these intersection points can be expressed in the form $\frac{a + b\sqrt{c}}{d}$, where $a$, $b$, $c$, and $d$ are positive integers such that $a$, $b$, $d$ have greatest common divisor equal to $1$, and $c$ is not divisible by the square of any prime. Find $a + b + c + d$.


 
 \textbf{Answer: }259



\textbf{Solution 1}

Note that the part of $f(x)$ in the first and fourth quadrants consists of lines of the form $y = x - 4k$ and $y = 4k + 2 - x$ for integers $k \geq 0$. Plugging in $x = 34y^2$ for $y = x - 4k$, we get $-34y^2 + y + 4k = 0,$ so the sum of the roots is $\frac{1}{34}$ by Vieta’s. Similarly, in the second equation, we get a sum of $-\frac{1}{34}.$ 

 
The top peaks of the graph form the line $y = 1$ and the bottom peaks form $y = -1.$ So the parabola will cross the positive line once it reaches the point $(34,1),$ and it will cross the negative line at $(34, -1).$For $y = x - 4k,$ we find $k = \frac{33}{4}$ for the positive $y$ values and $k = \frac{35}{4}$ for the negative y values. Both round down to $8,$ so the parabola will stop intersecting $y = x - 4k$ entirely after $k = 8.$

 
Now, for $y = 4k + 2 - x,$ we find $k = \frac{33}{4}$ for the positive y values and $k = \frac{31}{4}$ for the negative y values. $\frac{33}{4}$ rounds down to $8,$ but $\frac{31}{4}$ rounds down to $7,$ so the parabola will stop intersecting the negative $y$ values of $y = 4k + 2 - x$ after $k = 7,$ but it will intersect the positive $y$ values until $k = 8.$ 

 
There are $8$ values of $k$ from $0 - 7,$ so the sum contributed by $y = 4k + 2 - x$ until $k = 7$ is $-\frac{8}{34}.$ There are $9$ values of $k$ from $0 - 8,$ so the sum contributed by $y = x - 4k$ until $k = 8$ is $\frac{9}{34}.$ Adding these together we get $\frac{1}{34}.$ 

 
Now all that's left is to figure out the positive $y$ value at $k = 8$ for the line $y = 4k + 2 - x.$ We get 
 \[34-34y^2 = y.\] 
 \[y = \frac{-1 \pm \sqrt{68^2 + 1}}{68} = \frac{-1 + 5 \sqrt{185}}{68}.\] Adding our $\frac{1}{34}$ from earlier finally gives us $\frac{1 + 5 \sqrt{185}}{68} \implies \boxed{259}$.

 





\textbf{Solution 2}

Drawing the graph, we can use the sawtooth graph provided so nicely by MAA and draw out the parabola $x = 34y^2$. We realize that the sawtooth graph is just a bunch of lines where the positive slope lines are $y = x, y = x + 4, y = x + 8,...$. The intersections of these lines, along with the parabola are just solving the system of equations: $x = 34y^2$ and $y = x, y = x + 4, ...$. If we just take $y = x$ and $x = 34y^2$, we see that the sum of all $y$ by Vieta's is just $\frac{1}{34}$. Similarly, for $y = x + 4$, the sum of the roots by Vieta's is also $\frac{1}{34}$. So for all the positive slope lines intersecting with the parabola just gives the sum of all $y$ to continuously be $\frac{1}{34}$. Okay, now let's look at the negative slope lines. These will have equations of $y = 2 - x, y = 6 - x, y = 10 - x, ..., y = 34 - x, ...$. Similar to what we did above, we just set each of these equations along with the parabola $x = 34y^2$. The sum of all $y$ for each of these negative line intersections by Vieta's is $\frac{-1}{34}$. This keeps going for all of the lines until we reach $y = 34 - x$. Now, unfortunately, both solutions don't work as the negative solution is out of the range of $[1 , 3]$, $[5, 7]$ and so on. So we just need to take one solution for this and that being the positive one according to the graph. So we just need to solve $34 - y = 34y^2$ which means $34y^2 + y - 34 = 0$. Solving gives
 \[y = \frac{-1 \pm \sqrt{68^2 + 1}}{68} = \frac{-1 + 5 \sqrt{185}}{68}.\]So, the sums of the roots are $\frac{1}{34}$ + $\frac{-1}{34}$ + $\frac{1}{34}$ + .... + $\frac{-1}{34}$ + $\frac{1}{34}$ + $\frac{-1 + 5 \sqrt{185}}{68}.$ Nicely all the $\frac{1}{34}$ terms cancel out leaving with only one $\frac{1}{34}$ and $\frac{-1 + 5 \sqrt{185}}{68}.$ So the sum of these two is $\frac{1 + 5 \sqrt{185}}{68}.$ From there, the answer is $\boxed{259}$.

 \href{https://mathproblemsolvingskills.wordpress.com/}{MathProblemSolvingSkills.com}

 


 




	\item (\textit{AIME I 2025}) The set of points in $3$-dimensional coordinate space that lie in the plane $x+y+z=75$ whose coordinates satisfy the inequalities 
 \[x-yz<y-zx<z-xy\]forms three disjoint convex regions. Exactly one of those regions has finite area. The area of this finite region can be expressed in the form $a\sqrt{b},$ where $a$ and $b$ are positive integers and $b$ is not divisible by the square of any prime. Find $a+b.$


 
 \textbf{Answer: }510



\textbf{Solution 1}

Rewriting we have $z=75-x-y.$

 From the inequality $x-yz<y-zx$ we can rewrite to get,
 \[x-y(75-x-y)< y-x(75-x-y).\] 
 \[76x-76y+y^2-x^2<0.\] 
 \[(76-x-y)(x-y)<0.\]

 Similarly from the inequality $y-zx<z-xy$ we rewrite to get, 
 \[y-x(75-x-y)<(75-x-y)-xy.\]
 \[2y + 2xy + x^2 - 74x - 75 < 0.\]
 \[(x + 1)(2y + x - 75) < 0.\]Our next step is a visual which can be seen by roughly graphing the two inequalities. The first inequality is similar to a bow tie(you have to graph to see it lol) with bounds of $x-y=0$ and $76-x-y=0.$ The second one is a bow tie with edges of $x+1=0$ and $2y+x-75=0.$Here is the region of our solutions.
                    \begin{figure}[H]
                        \centering
                        \includegraphics[width=0.5\textwidth]{Solution/AIME_2025_I_12_Solution_Figure_1.png}
                    \end{figure}
                    It is simple to find the intersections of those which are $(-1,38,38),(25,25,25)$ and $(-1,-1,77).$ The sidelengths of this triangle are $39\sqrt{2},26\sqrt{6},13\sqrt{6}$ which is a $30-60-90$ gives us an area of 
 \[\frac{1}{2}(39\sqrt{2})(13\sqrt{6})=507\sqrt{3}\implies\boxed{510}.\]

 





\textbf{Solution 2}

Consider $x-yz<y-zx<z-xy$. From $x-yz<y-zx$, we find $z(y-x)>x-y$. Thus, if $x>y$, then $z<-1$, and if $x<y$, then $z>-1$. Similarly, taking another pair of the inequalities yields $y>-1$ when $z>x$ and $y<-1$ when $x>z$. Finally, taking the third pair yields $x>-1$ if $z>y$ and $x<-1$ if $z<y$.

 
Consider the first two resultant pairs of inequalities. Taking them pairwise (one from the first set and one from the second set) results in four cases:

 
1. Combining $z<-1$ if $x>y$ and $y>-1$ if $z>x$ yields $-1>z>x>y>-1$, a contradiction.

 2. Combining $z<-1$ if $x>y$ and $y<-1$ if $z<x$ yields $x,-1>y,z$.

 3. Combining $z>-1$ if $x<y$ and $y>-1$ if $z>x$ yields $y,z>-1,x$.

 4. Combining $z>-1$ if $x<y$ and $y<-1$ if $z<x$ yields $-1>y>x>z>-1$, a contradiction.

 
Now we have only two satisfactory inequalities. We now consider the third pair of inequalities ($x>-1$ if $z>y$ and $x<-1$ if $z<y$). Taking the two sets pairwise:

 
1. Combining $x,-1>y,z$ and $x>-1$ if $z>y$ yields $x>-1>z>y$. Consider some valid $x,y,z$ that satisfy $x+y+z=75$ and $x>-1>z>y$. We can infinitely increase $x$ while decreasing $y$ by the same amount, leading to another valid triple, so this case is infinite and we do not consider this case (for instance, if $x=100,y=-13,z=-12$, then $x=100+a,y=-13-a,z=-12$ is a valid triple for all nonnegative $a$).

 2. Combining $y,z>-1,x$ and $x>-1$ if $z>y$ yields $z>y>x>-1$. This case is finite due to the lower bound.

 3. Combining $x,-1>y,z$ and $x<-1$ if $z<y$ yields $-1>x>y>z$. There are no possible solutions since $x,y,z$ are negative from this inequality, but at least one must be positive to satisfy $x+y+z=75$, a contradiction.

 4. Combining $y,z>-1,x$ and $x<-1$ if $z<y$ yields $y>z>-1>x$. By the same argument as in Case 1, this is an infinite case.

 Thus we are tasked with finding the area of the figure formed by all triples $x,y,z$ satisfying $x+y+z=75$ and $z>y>x>-1$. We consider edge cases, so we maximize each variable by the largest amount possible to find three triples $(77,-1,-1),(38,38,-1),(25,25,25)$. We assume that these are the only edge cases (so the figure forms a triangle), and we can use the \href{https://artofproblemsolving.com/wiki/index.php?title=Distance\_formula}{Distance formula}. We find that the three side lengths of our triangle are $39\sqrt{2},13\sqrt{6},26\sqrt{6}$. These side lengths just so happen to form a $30-60-90$ triangle with legs $13\sqrt{6}$ and $39\sqrt{2}$, so the area of the triangle is

 
 \[\frac{1}{2}\cdot13\sqrt{6}\cdot39\sqrt{2}=507\sqrt{3}\]

 Thus the answer is $507+3=\boxed{510}$. ~eevee9406

 





\textbf{Solution 3}

Decomposing the inequality chain:
 \[x-yz<y-zx \quad \text{and} \quad y-zx<z-xy\]which is equivalent to
 \[(x-y)(z+1)<0 \quad \text{and} \quad (y-z)(x+1)<0\]Substituting $z$ with $z=75-x-y$ and simplifying yields
 \[(x-y)(x+y-76)>0 \quad \text{and} \quad (x+2y-75)(x+1)<0\]See that the solution to the first inequality is
 \[x-y>0, \, x+y-76>0 \quad \text{(I)} \quad \text{or} \quad x-y<0, \, x+y-76<0 \quad \text{(II)}\]Applying a similar method results in the solution to the second:
 \[x+2y-75>0, \, x+1<0 \quad \text{(III)} \quad \text{or} \quad x+2y-75<0, \, x+1>0 \quad \text{(IV)}\]Trying each grouping (i.e. let $\text{(I)}$ and $\text{(III)}$, $\text{(I)}$ and $\text{(IV)}$, $\text{(II)}$ and $\text{(III)}$, or $\text{(II)}$ and $\text{(IV)}$ be satisfied at the same time) and graphing shows that when $\text{(II)}$ and $\text{(IV)}$ are both satisfied, a triangle whose vertices are $(-1,38)$, $(-1,-1)$, and $(25,25)$ is formed. Further calculations show that the area of the triangle is $507$. However, this is not the final answer. We have projected the original shape to the $xy$-plane by substituting $z$. We know that for a surface defined by the equation $z=f(x,y)$, the area element $dS$ for this surface is given by
 \[dS=\sqrt{1+(f_x)^2+(f_y)^2}dxdy\]where $f_x$ and $f_y$ are the partial derivatives of the function $f(x,y)$ with respect to $x$ and $y$. For the plane $x+y+z=75$ where $f(x,y)=75-x-y$, computation gives
 \[f_x=-1, f_y=-1\]Substituting these into the original equation to get
 \[dS=\sqrt{3}dxdy\]This implies that to find the area of the original shape, we have to multiply the area of its projection on the $xy$-plane by $\sqrt{3}$. Therefore, the area of the original shape is $507\sqrt{3}$, with final answer $\boxed{510}$.

 \textbf{Supplementary Graph} 
                    \begin{figure}[H]
                        \centering
                        \includegraphics[width=0.5\textwidth]{Solution/AIME_2025_I_12_Solution_Figure_2.png}
                    \end{figure}
                    

 This is the real graph of the 3D space in coordinates. The area in BLUE is the space that satisfied the conditions, and it is 1/6 of a regular triangle.

 


 \textbf{Fakesolve: Fast but Flawed} Visualize or draw a $3$-dimensional graph of the following planes: $x+y+z=75$, $x=-1$, $y=-1$, $z=-1$, $x=y$, $y=z$, and $z=x$. The last $6$ planes (which be derived by solving for the zeroes of the inequalities $x-yz<y-zx$, $y-zx<z-xy$, $x-yz<z-xy$) form the plane $x+y+z=75$ into an triangle (which is equilateral by symmetry) with all $3$ of its medians cutting through it. 

 The vertices of the large equilateral triangle can be found by considering its smaller medial triangle formed by connecting the midpoints of the equilateral triangle. The coordinates of the medial triangle are formed by the intersections of the planes $x+y+z=75$, $x=y$, $z=-1$; $x+y+z=75$, $y=z$, $x=-1$; and $x+y+z=75$, $z=x$, $y=-1$. Solving these systems of equations gives the vertices $(38, 38, -1)$, $(-1, 38, 38)$, and $(38, -1, 38)$, respectively. Applying the Distance Formula, we get the side length of the medial triangle as $\sqrt{(x_2-x_1)^2+(y_2-y_1)^2+(z_2-z_1)^2}=\sqrt{(-1-38)^2+(38-38)^2+(38-(-1))^2}=39\sqrt{2}$.

 Then, the area of the original equilateral triangle is $4\cdot\frac{(39\sqrt{2})^2\cdot\sqrt{3}}{4}=3042\sqrt{3}$. Since the original equilateral triangle is divided into $6$ congruent $30-60-90$ right triangles, each of these triangles has area $3042\sqrt{3}/6=507\sqrt{3}$. Assuming that the desired convex region includes at least $1$ of these $30-60-90$ right triangles, the possible answers are $507\sqrt{3}, 1014\sqrt{3}, 1521\sqrt{3}, ...  3042\sqrt{3}$. However, if we compare each of these possible answers to the format $a\sqrt{b}$, the only one of these answers for which $a+b$ is less than $1000$ (so it fits into the AIME answer format of $3$ digits) is $507\sqrt{3}$, yielding $a+b=507+3=\boxed{510}$.

 \href{https://mathproblemsolvingskills.wordpress.com/}{MathProblemSolvingSkills.com}

 


 




	\item (\textit{AIME I 2025}) Alex divides a disc into four quadrants with two perpendicular diameters intersecting at the center of the disk. He draws $25$ more line segments through the disk, drawing each segment by selecting two points at random on the perimeter of the disk in different quadrants and connecting these two points. Find the expected number of regions into which these $27$ line segments divide the disk.


 
 \textbf{Answer: }204



\textbf{Solution 1}

First, we calculate the probability that two segments intersect each other. Let the quadrants be numbered $1$ through $4$ in the normal labeling of quadrants, let the two perpendicular diameters be labeled the $x$-axis and $y$-axis, and let the two segments be $A$ and $B.$

 $\textbf{Case 1:}$ Segment $A$ has endpoints in two opposite quadrants.

 
                    \begin{figure}[H]
                        \centering
                        \includegraphics[width=0.5\textwidth]{Solution/AIME_2025_I_13_Solution_Figure_1.png}
                    \end{figure}
                    

 This happens with probability $\frac{1}{3}.$ WLOG let the two quadrants be $1$ and $3.$ We do cases in which quadrants segment $B$ lies in.

 
                \begin{itemize}
                    \item Quadrants $1$ and $2,$ $2$ and $3,$ $3$ and $4,$ and $4$ and $1$: These share one quadrant with $A,$ and it is clear that for any of them to intersect $A,$ the endpoint that shares a quadrant with an endpoint of $A$ on a certain side of that endpoint of $A$ For example, if it was quadrants $1$ and $2,$ then the point in quadrant $1$ must be closer to the $x$-axis than the endpoint of $A$ in quadrant $1.$ This happens with probability $\frac{1}{2}.$ Additionally, segment $B$ has a $\frac{1}{6}$ to have endpoints in any set of two quadrants, so this case contributes to the total probability
                \end{itemize}


                
 \[\dfrac{1}{3}\left(\dfrac{1}{6}\cdot\dfrac{1}{2}+\dfrac{1}{6}\cdot\dfrac{1}{2}+\dfrac{1}{6}\cdot\dfrac{1}{2}+\dfrac{1}{6}\cdot\dfrac{1}{2}\right)=\dfrac{1}{9}\]

 
                \begin{itemize}
                    \item Quadrants $2$ and $4.$ This always intersects segment $A,$ so this case contributes to the total probability
                \end{itemize}


                
 \[\dfrac{1}{3}\cdot\dfrac{1}{6}=\dfrac{1}{18}\]

 
                \begin{itemize}
                    \item Quadrants $1$ and $3.$ We will first choose the endpoints, and then choose the segments from the endpoints. Let the endpoints of the segments in quadrant $1$ be $R_1$ and $R_2,$ and the endpoints of the segments in quadrant $3$ be $S_1$ and $S_2$ such that $R_1,R_2,S_1,$ and $S_2$ are in clockwise order. Note that the probability that $A$ and $B$ intersect is the probability that $A_1$ is paired with $B_1,$ which is $\dfrac{1}{2}.$ Thus, this case contributes to the total probability
                \end{itemize}


                
 \[\dfrac{1}{3}\cdot\dfrac{1}{6}\cdot\dfrac{1}{2}=\dfrac{1}{36}.\]

 $\textbf{Case 2:}$ Segment $A$ has endpoints in two adjacent quadrants.

 
                    \begin{figure}[H]
                        \centering
                        \includegraphics[width=0.5\textwidth]{Solution/AIME_2025_I_13_Solution_Figure_2.png}
                    \end{figure}
                    

 This happens with probability $\frac{2}{3}.$ WLOG let the two quadrants be $1$ and $2.$ We do cases in which quadrants segment $B$ lies in.

 
                \begin{itemize}
                    \item Quadrants $1$ and $3,$ $1$ and $4,$ $2$ and $3,$ and $2$ and $4.$ This is similar to our first case above, so this contributes to the total probability
                \end{itemize}


                
 \[\dfrac{2}{3}\left(\dfrac{1}{6}\cdot\dfrac{1}{2}+\dfrac{1}{6}\cdot\dfrac{1}{2}+\dfrac{1}{6}\cdot\dfrac{1}{2}+\dfrac{1}{6}\cdot\dfrac{1}{2}\right)=\dfrac{2}{9}\]

 
                \begin{itemize}
                    \item Quadrants $3$ and $4.$ This cannot intersect segment $A.$\item Quadrants $1$ and $2,$ Similar to our third case above, this intersects segment $A$ with probability $\frac{1}{2},$ so this case contributes to the total probability
                \end{itemize}


                
 \[\dfrac{2}{3}\cdot\dfrac{1}{6}\cdot\dfrac{1}{2}=\dfrac{1}{18}.\]Thus, the probability that two segments intersect is 
 \[\dfrac{1}{9}+\dfrac{1}{18}+\dfrac{1}{36}+\dfrac{2}{9}+\dfrac{1}{18}=\dfrac{17}{36}.\]Next, we will compute the expected number of intersections of a segment with the axes. WLOG let a segment have an endpoint in quadrant $1.$ Then, it will intersect each axis with probability $\dfrac{2}{3}$ because two out of the three remaining quadrants let it intersect a specific axis, so the expected number of axes a segment intersects is $\frac{4}{3}.$

 So, why do intersections matter? Because when adding a segment, it will pass through a number of regions, and for each region it passes through, it will split that region into two and add another region. The segment starts in a region, and for each intersection, it will enter another region, so the number of regions a segment passes through is $1$ more than the number of intersections with the axes and other segments. Thus, we have that by linearity of expectation the expected number of new regions created by adding a segment is
 \[\dfrac{17}{36}\cdot(\text{number of segments already added})+\dfrac{4}{3}+1,\]so the number of new regions added in total by $25$ segments again by linearity of expectation is
 \[\sum_{k=0}^{24}\left(\dfrac{17}{36}k+\dfrac{7}{3}\right)=\dfrac{17}{36}\cdot \dfrac{24\cdot 25}{2}+\dfrac{25\cdot 7}{3}\]which simplifies to $200$ as the expected number of new regions added by the $25$ segments. The axes create $4$ regions to begin with, so our answer is
 \[200+4=\boxed{204}.\]

 





\textbf{Solution 2}

Claim: If N line segments are drawn in the disk so that no three meet at an interior point (true with probability 1 here), thenR = 1 + N + I,where I is the total number of interior intersection points among the segments.

 Reason: When you draw a new segment that meets previous segments in k interior points, it is cut into k+1 pieces, creating k+1 new regions. Summing “+1” over all segments gives +N, and summing k over all segments counts each intersection once, giving +I.

 Here N = 27, soE[R] = 1 + 27 + E[I] = 28 + E[I].

 Step 1: Intersections involving the two perpendicular diameters.The two diameters intersect once at the center.

 Now consider a random chord: endpoints are chosen uniformly on the circle, conditioned to lie in different quadrants. Thus the unordered quadrant-pair is uniform among the C(4,2)=6 pairs.

 Fix one diameter (say the horizontal). A chord crosses it iff its endpoints lie in opposite half-disks, i.e. it uses one quadrant above and one below. Among the 6 quadrant-pairs, 4 have this property, soP(chord intersects a fixed diameter) = 4/6 = 2/3.

 There are 25 chords and 2 diameters, soE[(diameter–chord intersections)] = 25 · 2 · (2/3) = 100/3.

 So far:E[I] = 1 + 100/3 + E[(chord–chord intersections)].

 Step 2: Probability that two random chords intersect.Label the four quadrant arcs cyclically A, B, C, D. A chord type is an unordered pair of distinct arcs; there are 6 types:AB, BC, CD, DA (adjacent) and AC, BD (opposite).

 Because all arcs have equal length, conditioning on “different quadrants” makes the 6 chord types equally likely; hence an ordered pair of chord types is uniform among 6² = 36 cases.

 For each ordered type-pair:• If the two chords are AB and CD (or CD and AB), or BC and DA (or DA and BC), they cannot intersect (their endpoints lie in two disjoint arcs in a non-alternating way): probability 0. That gives 4 ordered cases.• If the two chords are AC and BD (or BD and AC), they always intersect: probability 1. That gives 2 ordered cases.• In every other ordered case, the chords share exactly one arc (e.g. AB with AC shares A, etc.). Then, whether they intersect depends only on which of the two points on that shared arc comes first along the circle, which is equally likely. So the intersection probability is 1/2. There are 36 − 4 − 2 = 30 such ordered cases.

 ThereforeP(two random chords intersect) = [2·1 + 30·(1/2) + 4·0] / 36= (2 + 15)/36= 17/36.

 There are C(25,2) = 300 chord pairs, soE[(chord–chord intersections)] = 300 · (17/36) = 425/3.

 Step 3: Assemble.E[I] = 1 + 100/3 + 425/3= 1 + 525/3= 1 + 175= 176.

 ThusE[R] = 28 + 176 = 204.

 Answer: 204.

 \href{https://mathproblemsolvingskills.wordpress.com/}{MathProblemSolvingSkills.com}

 


 




	\item (\textit{AIME I 2025}) Let $ABCDE$ be a convex pentagon with $AB=14,$ $BC=7,$ $CD=24,$ $DE=13,$ $EA=26,$ and $\angle B=\angle E=60^{\circ}.$ For each point $X$ in the plane, define $f(X)=AX+BX+CX+DX+EX.$ The least possible value of $f(X)$ can be expressed as $m+n\sqrt{p},$ where $m$ and $n$ are positive integers and $p$ is not divisible by the square of any prime. Find $m+n+p.$


 
 \textbf{Answer: }060



\textbf{Solution 1}

Assume $AX=a$, $BX=b$, $CX=c$, $DX=d$, $EX=e$, by Ptolemy inequality we have $a+2d\geq \sqrt{3}XE; a+2c\geq \sqrt{3}BX$, and as we are minimizing the value of the LHS, we want the inequality reached; this happens when both quadrilaterals $CXAB$ and $AXDE$ are concyclic. We notice that we have right angles $\angle{BXA}=\angle{BCA}$ and $\angle{EDA}=\angle{EXA}$ as they both subtend a diameter of circles by noticing that we have formed Pythagorean triples (along with being told that $\angle{B}=\angle{E}=60^{\circ}$. )And so, since we have supplementary angles, $B$, $X$, and $E$ are collinear. Thus, we add up the two equalities to get $2a+2c+2d=\sqrt{3}(XE+BX)$, or $a+c+d=\frac{\sqrt{3}}{2}(BE)$. Thus we want to find $a+b+c+d+e=(1+\frac{\sqrt{3}}{2})BE$.

 By Law of Cosines, we know that $\cos(\angle{DAC})=\frac{1}{7}$, and since $\angle{CAB}=\angle{DAE}=30^{\circ}$, we know $\cos(\angle{EAB})=cos(\angle{CAB}+\angle{DAE}+\arccos(\frac{1}{7}))=cos(60^{\circ}+\arccos(\frac{1}{7}))$, which we find with the cosine angle addition formula and the fact that $\sin(\arccos(\frac{1}{7}))=\sqrt{1-(\frac{1}{7})^2}$. Then, once again with Law of Cosines, we find  that $BE=38$, and thus the minimum of $a+b+c+d+e$, or $(1+\frac{\sqrt{3}}{2})BE$, is $38+19\sqrt{3}\implies\boxed{060}$.

 Rewritten and edited by \href{https://artofproblemsolving.com/wiki/index.php?title=User:Juwushu}{ Juwushu}.

 





\textbf{Solution 2}


                    \begin{figure}[H]
                        \centering
                        \includegraphics[width=0.5\textwidth]{Solution/AIME_2025_I_14_Solution_Figure_1.png}
                    \end{figure}
                    Firstly, note that $\triangle ABC$ and $\triangle ADE$ are just 30-60-90 triangles. Let $X$ be the Fermat point of $\triangle ACD$, with motivation stemming from considering the pentagon as $\triangle ACD$ with the two 30-60-90 extensions. Note that $AX+CX+DX$ is minimized at this point when $\angle AXC=\angle CXD=\angle AXD=120^{\circ}$. Because we have $\angle ABC=\angle AED=60^{\circ}$, then $ABCX$ and $AXDE$ are both cyclic. Then we have $\angle AXE=\angle ADE=90^{\circ}$ and $\angle BXA=\angle BCA=90^{\circ}$. Then it turns out that we actually have $\angle BXE=90^{\circ}+90^{\circ}=180^{\circ}$, implying that $B$, $X$ and $E$ are collinear. Now, by the triangle inequality, we must have $BX+XE\geq BE$, with equality occurring when $X$ is on $BE$. Thus $AX+CX+DX$ and $BX+EX$ are minimized, so this point $X$ is our desired point.

 Firstly, we will find $BX+EX=BE$. We have that $AC=7\sqrt{3}$ and $AD=13\sqrt{3}$, so applying the Law of Cosines in $\triangle ACD$, we get 
 \[147+507-2(7\sqrt{3})(13\sqrt{3})\cos (\angle CAD)=576\implies \cos(\angle CAD)=\frac{1}{7}.\] It follows as a result that $\sin (\angle CAD)=\frac{4\sqrt{3}}{7}$. Then we want to find $\cos (\angle BAE)$. We can do this by seeing 
 \[\cos (\angle BAE)=\cos (\angle CAD+60^{\circ})=\cos (\angle CAD)\cos 60^{\circ}-\sin (\angle CAD)\sin 60^{\circ}=\frac{1}{7}\cdot \frac{1}{2}-\frac{4\sqrt{3}}{7}\cdot \frac{\sqrt{3}}{2}=-\frac{11}{14}.\] Applying the Law of Cosines again in $\triangle BAE$, then because $AB=14$ and $AE=26$, we have 
 \[14^2+26^2-2(14)(26)\left (-\frac{11}{14}\right )=196+676-2\cdot 26\cdot (-11)=872+572=1444=BE^2,\] so it follows that $BE=38=BX+EX$.

 Now, we will find the value of $AX+CX+DX$. Construct a point $P$ outside such that $\triangle CPD$ is equilateral, as shown. By property of fermat point, then $A$, $X$, and $P$ are collinear. Additionally, $\angle CXD=120^{\circ}$, so $CPDX$ is cyclic. Applying Ptolemy's Theorem, we have that $(CX)(PD)+(CP)(XD)=(XP)(CD)$. But since $\triangle CPD$ is equilateral, it follows that $CX+DX=PX$. Then $AX+CX+DX=AX+PX=AP$, so we wish to find $AP$. Applying the Law of Cosines in $\triangle ACD$, we have that 
 \[(13\sqrt{3})^2+24^2-2(13\sqrt{3})(24)\cos (\angle ADC)=(7\sqrt{3})^2\implies \cos (\angle ADC)=\frac{\sqrt{3}}{2}\implies \angle ADC=30^{\circ}.\] Then because $\angle CDP=60^{\circ}$, then $\angle ADP=90^{\circ}$, so we can find $AP$ simply with the Pythagorean Theorem. We know $AD=13\sqrt{3}$ and $DP=CD=24$, so $AP=\sqrt{(13\sqrt{3})^2+24^2}=19\sqrt{3}$.

 We then have $f(X)=AX+BX+CX+DX+EX=(BX+EX)+(AX+CX+DX)=BE+AP=38+19\sqrt{3}$, which is our minimum value. Therefore, the answer to the problem is $38+19+3=\boxed{060}$.

 





\textbf{Solution 3(Fermat Point)}

$AE = 26$, $DE = 13$, and $\angle{E}={60}^{\circ}$, which means that $\triangle{AED}$ is a 30-60-90 triangle, so $AD = 13\sqrt3$, $\angle{EAD}={30}^{\circ}$, and $\angle{ADE}={90}^{\circ}$. Similar with $\triangle{ABC}$, $AD = 7\sqrt3$, $\angle{BAC}={30}^{\circ}$, and $\angle{ACB}={90}^{\circ}$

 To solve the question, we would have to locate point $X$ first. We first consider the points $B$ and $E$. For the distance of $X$ to $B$ and $E$ to become the shortest, $X$ should lay on $BE$. For $X$ to be closest to point $A$, it should be on the foot of perpendicular from $A$ to line $BE$. 

 Consider about $C$ and $D$. $\angle{ADE}=\angle{AXE}={90}^{\circ}$, so $AXDE$ is cyclic. Therefore, $\angle{EXD}=\angle{EAD}=\angle{BXC}={30}^{\circ}$. $\angle{DXC}=\angle{AXD}=\angle{AXD}={120}^{\circ}$, so $X$ is coincidently the Fermat Point of $ADC$.

 To calculate the $f(X)$, we divide it into 2 parts: the sum of distance to $A$, $C$, and $D$ and the sum of distance to $B$ and $E$. Applying Ptolemy's Theorem in $AXDE$ and $AXCB$,
 \[AX+2DX=\sqrt3EX\]
 \[and\]
 \[AX+2CX=\sqrt3BX\]We get that $AX+DX+CX = \frac{\sqrt3}{2}\cdot(EX+BX) = \frac{\sqrt3}{2}BE$

 Construct equilateral triangle $\triangle{ADE}$ outside of $\triangle{ADC}$ on side $AD$. Because $X$ is the Fermat Point, $FC=AX+DX+CX$. To calculate $FC$, we needed to utilize $\angle{FDC}$

 
 \[\angle{FDC}=\angle{FDA}+\angle{ADC}={60}^{\circ}+\angle{ADC}\].

 From $\triangle{ADC}$, we know:\begin{align*}\cos{\angle{ADC}} & = \frac{(13\sqrt3)^2+{24}^2-(7\sqrt3)^2}{2\cdot13\sqrt3\cdot24}\\&= \frac{936}{2\cdot13\sqrt3\cdot24}\\&= \frac{\sqrt3}{2}\\\end{align*}


 This shows that $\angle{ADC} = {30}^{\circ}$, which means that $\angle{FDC} = {60}^{\circ}+{30}^{\circ}={90}^{\circ}$

 Using the Pythagorean Theorem,
 \[FC = \sqrt{(13\sqrt3)^2+24^2)}=19\sqrt3=AX+DX+CD\]
 \[f(X) = 19\sqrt3(1+\frac{2}{\sqrt3})=38+19\sqrt3\] 

 The answer is $38+19+3=\boxed{060}$

 \href{https://mathproblemsolvingskills.wordpress.com/}{MathProblemSolvingSkills.com}

 Great solution!

 




	\item (\textit{AIME I 2025}) Let $N$ denote the number of ordered triples of positive integers $(a, b, c)$ such that $a, b, c \leq 3^6$ and $a^3 + b^3 + c^3$ is a multiple of $3^7$. Find the remainder when $N$ is divided by $1000$.


 
 \textbf{Answer: }735



\textbf{Solution 1}

First, state the LTE Lemma for $p = 3, n = 3$, which we might use.

 


 \begin{verbatim}
$\bullet$ $\nu_3(n) = \begin{cases}         \max \{k : 3^k \mid n\} &n\neq 0\\         \infty &n=0     \end{cases}$
\end{verbatim}\begin{verbatim}
$\bullet$ If $3 \nmid x, 3\nmid y, 3\mid x+y$, then $\nu_3(x^3+y^3) = \nu_3(x+y) + \nu_3(3) = \nu_3(x+y) + 1$      $\bullet$ If $3 \nmid x, 3\nmid y, 3\mid x-y$, then $\nu_3(x^3-y^3) = \nu_3(x-y) + \nu_3(3) = \nu_3(x-y) + 1$   
\end{verbatim}Then we classify all cube numbers $\mod{3^7}$

 


 \begin{verbatim}
$\bullet$ $A = \{a : 1\leq a \leq 3^6, 9\mid a\}$
\end{verbatim}\begin{verbatim}
We can write $A = \{a: a=9k, 1\leq k \leq 3^4\}$, so $|A| = 3^4$.          $\bullet$ If $k \equiv 0 \mod{3}$, $k^3 \equiv 0 \mod{3}, a^3 \equiv 3^6k^3 \equiv 0 \mod{3^7}$, there are 27 roots.       $\bullet$ If $k \equiv 1 \mod{3}$, $k^3 \equiv 1 \mod{3}, a^3 \equiv 3^6k^3 \equiv 3^6 \mod{3^7}$, there are 27 roots.       $\bullet$ If $k \equiv 2 \mod{3}$, $k^3 \equiv 2 \mod{3}, a^3 \equiv 3^6k^3 \equiv 2\times 3^6 \mod{3^7}$, there are 27 roots.         $\bullet$ $B = \{a : 1\leq a \leq 3^6, 3\mid a, 9 \nmid a\}$
\end{verbatim}\begin{verbatim}
We can write $B = \{a: a=9k+3 \text{ or }a=9k+6, 0\leq k < 3^4\}$, so $|B| = 2 \times 3^4$. 
\end{verbatim}\begin{verbatim}
For $x,y \in \{a: a=9k+3, 0\leq k < 3^4\}$, let $x = 3a, y= 3b$ and hence $3 \nmid a, 3\nmid b, 3\mid a-b$.          $(3a)^3 - (3b)^3 \equiv 0 \mod{3^7}$       $\iff a^3 - b^3 \equiv 0 \mod{3^4}$       $\iff \nu(a^3-b^3) \geq 4$       $\iff \nu(a-b) \geq 3$       $\iff a - b\equiv 0 \mod{3^3}$       $\iff x - y\equiv 0 \mod{3^4}$      For $k = 0, ..., 8$, each has 9 roots.      Since $(9k+3)^3 \equiv 3^6k^3+3^6k^2+3^5k + 3^3 \equiv 3^5m + 3^3 \mod{3^7}$, $0\leq m \leq 8$. They are the corresponding classes.
\end{verbatim}\begin{verbatim}
Same apply to $x,y \in \{a: a=9k+6, 0\leq k < 3^4\}$. For $k = 0, ..., 8$, each has 9 roots.
\end{verbatim}\begin{verbatim}
Since $(9k+6)^3 \equiv 3^6k^3+2\times3^6k^2+4\times3^5k + 8\times3^3 \equiv 3^5m - 3^3 \mod{3^7}$, $0\leq m \leq 8$. They are the corresponding classes.
\end{verbatim}\begin{verbatim}
$\bullet$ $C = \{a : 1\leq a \leq 3^6, 3\nmid a\}$
\end{verbatim}\begin{verbatim}
We write $C = \{a: a=3k+1 \text{ or }a=3k+2, 0\leq k < 3^5\}$, so $|C| = 2 \times 3^5$.
\end{verbatim}\begin{verbatim}
For $a,b \in \{a: a=3k+1, 0\leq k < 3^5\}$, then $3 \nmid a, 3\nmid b, 3\mid a-b$.          $a^3 - b^3 \equiv 0 \mod{3^7}$       $\iff \nu(a^3-b^3) \geq 7$       $\iff \nu(a-b) \geq 6$       $\iff a - b\equiv 0 \mod{3^6}$   For $k = 0, ..., 3^5-1$, 1 root each.      $(3k+1)^3 \equiv 3^3k^3+3^3k^2+3^2k + 1 \equiv 3^2m + 1 \mod{3^7}$, $0\leq m < 3^5$. They are the corresponding classes.
\end{verbatim}\begin{verbatim}
Same apply to $x,y \in \{a: a=3k+2, 0\leq k < 3^5\}$. For $k = 0, ..., 3^5-1$, 1 root each.
\end{verbatim}\begin{verbatim}
$(3k+2)^3 \equiv 3^3k^3+2\times3^3k^2+4\times3^2k + 8 \equiv 3^2m - 1 \mod{3^7}$, $0\leq m < 3^5$. They are the corresponding classes.
\end{verbatim}
Summarized the results:

 \begin{verbatim}
$\bullet$ $A$: If $x \equiv 0, 3^6, 2\times3^6 \mod{3^7}$, then $x$ has 27 roots. $|A| = 3^4$.   $\bullet$ $B$: If $x \equiv 3^5m \pm 3^3 \mod{3^7}$, then $x$ has 9 roots. $|B| = 2\times3^4$.   $\bullet$ $C$: If $x \equiv 3^2m \pm 1 \mod{3^7}$, then $x$ has 1 root. $|C| = 2\times3^5$.   $\bullet$ Otherwise, $x$ has no roots.
\end{verbatim}
We do the final combinatorial problem.

 


 \begin{verbatim}
$\bullet$ Case: $A,A,A$: $|A| \times |A| \times 27 = \boxed{3 \times 3^{10}}$   $\bullet$ Case $A,A,B$: No solution.   $\bullet$ Case $A,A,C$: No solution.   $\bullet$ Case: $A,B,B$: $3 \times |A| \times |B| \times 9 = \boxed{6 \times 3^{10}}$   $\bullet$ Case $A,A,B$: No solution.   $\bullet$ Case: $A,C,C$: $3 \times |A| \times |C| \times 1 = \boxed{2 \times 3^{10}}$   $\bullet$ Case $B,B,C$: No solution.   $\bullet$ Case $B,B,B$: No solution.   $\bullet$ Case $B,C,C$: $3 \times |B| \times |C| \times 1 = \boxed{4 \times 3^{10}}$   $\bullet$ Case $C,C,C$: No solution.    
\end{verbatim}Total is $(3+6+2+4)3^{10}=15\times 3^{10}$.

 $3^5 = 243 \equiv 43 \mod{200},43^2=1600+240+9\equiv49\mod{200}$

 Hence $15\times3^{10}\equiv15\times49\equiv735\mod{1000}$

 Answer is $\boxed{735}$.

 





\textbf{Solution 2}

We are given that$N = \#\{(a,b,c) \in [1,3^6]^3 : 3^7 \mid (a^3+b^3+c^3)\},$and we wish to compute the remainder when \(N\) is divided by 1000.

 Since \(3^6=729\) and \(3^7=2187\), the variables \(a\), \(b\), and \(c\) range over the set \(\{1,2,\dots,729\}\).

 A standard approach is to use exponential sums to detect the divisibility condition. For any integer \(n\) we have the identity$\frac{1}{3^7}\sum_{t=0}^{3^7-1} e^{\frac{2\pi i\,t\,n}{3^7}} = \begin{cases} 1, & \text{if } 3^7\mid n,\\[1mm] 0, & \text{otherwise.} \end{cases}$
 \[\]Thus, we can write$N = \frac{1}{3^7}\sum_{t=0}^{3^7-1} \sum_{a=1}^{3^6} e^{\frac{2\pi i\,t\,a^3}{3^7}} \sum_{b=1}^{3^6} e^{\frac{2\pi i\,t\,b^3}{3^7}} \sum_{c=1}^{3^6} e^{\frac{2\pi i\,t\,c^3}{3^7}}.$Define$S(t)=\sum_{a=1}^{3^6} e^{\frac{2\pi i\,t\,a^3}{3^7}}.$Then,$N=\frac{1}{3^7}\sum_{t=0}^{3^7-1} \bigl(S(t)\bigr)^3.$

 $\medskip$

 Notice that when \(t=0\) we have$S(0)=\sum_{a=1}^{3^6} 1=3^6.$Thus, the \(t=0\) term contributes$\frac{1}{3^7}(3^6)^3 = \frac{3^{18}}{3^7}=3^{11}.$

 It turns out that for \(t\not\equiv 0\pmod{3^7}\), one can show (using properties of cubic residues modulo \(3^7\)) that$S(t)=0.$ [\textbf{Note: I don't believe this is true.  Try it manually for the modulo $3^2$ case.}]

 Therefore, only the \(t=0\) term contributes to the sum, and we obtain$N=3^{11}.$

 Since$3^{11}=177147,$we compute$(5 \cdot 177147) \mod 1000 = 735.$ [\textit{Note: The multiplication by 5 isn't justified. The computed $N$ above is simply wrong by a factor of 5. Unfortunately, I don't know how to fix this approach.}](context: he chat gptd the solution like recently after the test was finished so we didnt know what the real answer was. i bashed with a simple program and determined his solution was off by a factor of 5 so he just did that at the end ;-;)~shreyan.chethan

 





\textbf{Solution 3}

For $3^{7} \mid a^{3} + b^{3} + c^{3}$ to happen, we need $3 \mid a^{3} + b^{3} + c^{3}$. So we can define classes:

 $a,b,c \equiv 0\mod 3 \hspace{1 cm} \textbf{(1)}$

 $a,b,c \equiv 1\mod 3 \hspace{1 cm} \textbf{(2)}$

 $a,b,c \equiv 2\mod 3 \hspace{1 cm} \textbf{(3)}$

 $a \equiv 0\mod 3, b \equiv 1\mod 3, c \equiv 2\mod 3 \hspace{1 cm} \textbf{(4)}$

 $a \equiv 0\mod 3, b \equiv 2\mod 3, c \equiv 1\mod 3 \hspace{1 cm} \textbf{(5)}$

 $a \equiv 1\mod 3, b \equiv 0\mod 3, c \equiv 2\mod 3 \hspace{1 cm} \textbf{(6)}$

 $a \equiv 1\mod 3, b \equiv 2\mod, c \equiv 0\mod 3 \hspace{1 cm} \textbf{(7)}$

 $a \equiv 2\mod 3, b \equiv 0\mod 3, c \equiv 1\mod 3 \hspace{1 cm} \textbf{(8)}$

 $a \equiv 2\mod 3, b \equiv 1\mod 3, c \equiv 0\mod 3 \hspace{1 cm} \textbf{(9)}$

 If we consider class (1), we have $a = 3a'$ and the others homogeneously. Then $a', b', c' \leq 3^{5}$ and $3^{7} \mid 3^{3}(a'^{3} + b'^{3} + c'^{3})$. So $3^{4} \mid a'^{3} + b'^{3} + c'^{3}$. Then we need to choose another class so we choose class (1) again. Then $a = 3a''$ and the other homogeneously. Then $a'', b'', c'' \leq 3^{4}$ and $3 \mid a''^{3} + b''^{3} + c''^{3}$. If we choose class (1) again, we would have satisfied the inequality and it would just be a number of multiples of $3$ between $1$ and $81$ so $27^{3} = 3^{9}$ triples. Then we choose class (2) to get the number of numbers that are 1(mod 3) between $1$ and $81$ which is also $27^{3} = 3^{9}$. In fact, we would just have $3^{9}$ triples for each of the classes so $9(3^{9}) = 3^{11}$ triples for this case for (1)-(1) classes. 

 Now, if we choose class (1) again, then we would have $3^{4} \mid a'^{3} + b'^{3} + c'^{3}$ and $a', b', c' \leq 3^{5}$. But if we choose class (2) this time, then we would have $a' = 3a'' + 1$ and the others. Then $a'', b'', c'' \leq 80$. If we expand and simplify, we would have $3^{3}$ divides a number that is 1(mod 3) which is clearly impossible. Same thing for (1)-(3) class arrangement. So we can skip to (1)-(4) class. 

 So we start with $3^{4} \mid a'^{3} + b'^{3} + c'^{3}$ and $a', b', c' \leq 3^{5}$. Choosing class (4) gives $a' = 3a'', b' = 3b'' + 1, c' = 3c'' - 1$. Expanding and simplifying we would need $3^{2} \mid a''^{3} + b''^{3} + c''^{3} + b''^{2} - c''^{2} + b'' + c''$. Now, we would just need to do all $a'', b'', c''$ residues modular 3 such that 3 divides the expression. We can classify them as:

 $a,b,c \equiv 0\mod 3$

 $a,c \equiv 0\mod 3, b \equiv 1\mod 3$

 $a \equiv 0\mod 3, b \equiv 1\mod 3, c \equiv 2\mod 3$

 $a \equiv 0\mod 3, b \equiv 2\mod 3, c \equiv 1\mod 3$

 $a,b \equiv 0\mod 3, c \equiv 2\mod 3$

 $a \equiv 1\mod 3, b \equiv 2\mod 3, c \equiv 0\mod 3$

 $a \equiv 1\mod 3, b,c \equiv 2\mod 3$

 $a \equiv 2\mod 3, b \equiv 0\mod 3, c \equiv 1\mod 3$

 $a \equiv 2\mod 3, b,c \equiv 1\mod 3$

 If we consider the first case, then we have $a'' = 3a'''$ and the others. Notice $1 \leq a''' \leq 27, 0 \leq b''' \leq 26, 1 \leq c''' \leq 27$. If we simplify with this knowledge and plug it into the original expression, we just need $3 \mid b''' + c'''$. Then, we would have $b''', c''' \equiv 0\mod 3$ to get $9(9) = 3^{4}$ solutions. We can keep going for $b''' \equiv 1\mod 3$ and $c''' \equiv 2\mod 3$ to get another $9(9) = 3^{4}$ solutions. In total, we have $3(3^{4}) = 3^{5}$ solutions after considering all $b'''$ and $c'''$ residues mod 3. But the value of $a'''$ doesn't matter so that generates another $27 = 3^{3}$ solutions. In total, we have $3^{3} \cdot 3^{5} = 3^{8}$ solutions. Now, we move on to the next case. Choosing this case will give that $a'' = 3a''', b'' = 3b''' + 1, c'' = 3c'''$ where $1 \leq a''' \leq 27, 0 \leq b''' \leq 26, 1 \leq c''' \leq 27$. So we need $3^{2} \mid 3^{3} a'''^{3} + 3^{3} b'''^{3} + 3^{3} b'''^{2} + 3^{2} b''' + 1 + 3^{3} c'''^{3} + 3^{2} b'''^{2} + 2 \cdot 3 \cdot b''' + 1 - 3^{2} c'''^{2} + 3b''' + 1 + 3c'''$. Then simplifying gives $3^{2} \mid 2 \cdot 3 \cdot b''' + 3 + 3b''' + c'''$. This yields $3 \mid 2b''' + 1 + b''' + c'''$ which gives $3 \mid c''' + 1$ so $c''' \equiv 2\mod 3$. So $a'''$ and $b'''$ don't really matter for this case so we have $27(27) = 3^{6}$ possible pairs for $a''', b'''$. Then we need to count the number of numbers that are 2(mod 3) between 1 and 27 inclusive which is 9. So we have $9(3^{6}) = 3^{8}$ total triples for case 2. The next case yields $3^{9}$ solutions following similar logic. In fact, all the remaining classes each give $3^{8}$ solutions for a total of $8(3^{8}) + 3^{9} = 11 \cdot 3^{8}$. But this is just for the class (1)-(4). Since classes (5),...,(9) are the same thing except it is just permutations of class (4), we will still generate the same number of solutions as (1)-(4) arrangement. So we would have $6$ classes each giving $11(3^{8})$ solutions for a total of $22(3^{9})$ solutions.

 So we have completed all class (1) cases. Now we proceed to class (2). So if we start with class (2), then we have $a = 3a' + 1$ and homogeneously the others. Then, we have $3^{6} \mid 3^{2}(a'^{3} + b'^{3} + c'^{3}) + 3^{2}(a'^{2} + b'^{2} + c'^{2}) + 3(a' + b' + c') + 1$ which is clearly impossible the $RHS$ expression is 1(mod 3). So we cannot begin with class (2). Similarly, we cannot begin with class (3). So we must proceed to class (4).

 Starting with class (4) gives us that $a = 3a', b = 3b' + 1, c = 3c' - 1$. So $a', c' \leq 3^{5}$ and $b' \leq 242$. We would get $3^{5} \mid 3(a'^{3} + b'^{3} + c'^{3}) + 3(b'^{2} - c'^{2}) + b' + c'$. So we need $3 \mid b' + c'$ at first. Following our logic by building on residues mod 3, we have:

 $b', c' \equiv 0\mod 3$

 $b' \equiv 1\mod 3, c' \equiv 2\mod 3$

 $b' \equiv 2\mod 3, c' \equiv 1\mod 3$

 Following case 1, we have $b' = 3b'', c' = 3c''$. This part in essence is basically all the casework bash we did for class (1) arrangements except this time, the powers of 3 are higher. We have to substitute and consider classes mod 3 yet again(I won't go through all the steps here because it's pretty much similar to how we found the above cases). We end with $3^{9}$ total cases. This would happen for all classes (4),...,(9) so we have to multiply this by 6 to get $6 \cdot 3^{9}$ cases. The casework we did here was based on (4)-(5) class arrangements, I didn't actually consider (4)-{(1),(2),(3)} class arrangements. Doing the casework for all those classes, we end up with $8(3^{9})$ cases(I won't repeat the steps here because it's practically the same thing, bashing all modular residues mod 3. Essentially we would get another copy of the $6(3^{9})$ and then doing the casework for (4)-(2) and (4)-(3), we would get $6(3^{8})$ options for a total of $8(3^{9})$ triples). Thus, we actually have $6(3^{9}) + 8(3^{9}) = 14(3^{9})$ total cases for these subclasses. And finally, we are ready to add these all up.

 We have $3^{11}$ total from class (1) trivial arrangements. Then, we have $22(3^{9})$ total cases from the (1)-(4) arrangements e.t.c. And finally, we have $14(3^{9})$ cases from the (4)-{(5)...(9)} classes. We have $3^{11} + 22(3^{9}) + 14(3^{9}) = 3^{11} + 36(3^{9}) = 3^{11} + 4(3^{11}) = 5(3^{11})$. We can calculate this mod 1000 by simply noting $3^{10} \equiv 49\mod 1000$ because $3^{10} \equiv 243^{2}\mod 1000$ which tells us that the answer is just $5 \cdot 3 \cdot 49 = \boxed{735}$.

 





\textbf{Solution 4(Recursion)}

Let $N(n, k)$ equal to number of triples in the range $(1, 2, ..., 3^k)$ for which $3^n | a^3 + b^3 + c^3$. Define $f(n,k)$ to be the number of triples such that $3 | a,b,c$ and define $g(n,k)$ to be the number of triples otherwise. It's obvious that $N(n, k) = f(n,k) + g(n,k)$. We can actually find a nice recursive definition for $f(n,k)$. We will replace $a = 3a_1, b = 3b_1, c = 3c_1 \implies 3^n | 3^3 a_1^3 + 3^3 b_1^3 + 3^3 c_1^3 \implies 3^{n - 3} | a_1^3 + b_1^3 + c_1^3$ which implies that the number of triples for this is simply $N(n - 3, k - 1)$ since we lower the range by one exponent of $3$ since we're dividing all of $a,b,c$ by $3$. Finding $g(n,k)$ is a bit harder. Note that all the residues modulo $3$ that give $3 | a^3 + b^3 + c^3$ are $0, 1, 2$ modulo $3$ and $1,1,1$ modulo $3$ and $2,2,2$ modulo $3$. Note that if $a, b, c \equiv 1\mod 3$, then

 $3^n | (3a_1 + 1)^3 + (3b_1 + 1)^3 + (3c_1 + 1)^3 = 3^3(a_1^3 + b_1^3 + c_1^3) + 3^3(a_1^2 + b_1^2 + c_1^2) + 3^2(a_1 + b_1 + c_1) + 3 \implies 3^{n - 1} | 3^2(a_1^3 + b_1^3 + c_1^3) + 3^2(a_1^2 + b_1^2 + c_1^2) + 3(a_1 + b_1 + c_1) + 1$ which is impossible since the RHS is $1$ modulo $3$. So the only residues that can work are $a,b,c$ are congruent to $0,1,2$ modulo $3$ with there being $3! = 6$ permutations. For now, we'll assume WLOG that $a \equiv 0\mod 3, b \equiv 1\mod 3, c \equiv 2\mod 3$ and then multiply by $6$ at the end. By Hensel's Lemma, we have exactly one solution for $c \equiv -(a^3 + b^3)\mod 3^n$ in modulo $3^{n - 1}$ so we can just count the number of triples.

 For $a$ since $3 | a$, we have $\frac{3^k}{3} = 3^{k - 1}$ choices for $a$. The same thing for $b$ with $3^{k - 1}$ choices. Since we're in modulo $3^{n - 1}$ for $c$, we must have

 $\frac{3^k}{3^{n - 1}} = 3^{k - n + 1}$.

 Finally, we can evaluate $g(n, k)$

 $g(n, k) = 6 \cdot 3^{k - 1} \cdot 3^{k - 1} \cdot 3^{k - n + 1} = 2 \cdot 3^{3k - n}$.

 Note that we found a recursive definition for $f(n, k) = N(n - 3, k - 1)$ no if we plug in $n = 7, k = 6$ we have

 $f(7, 6) = N(4, 5) = f(4, 5) + g(4, 5) = N(1, 4) + g(4, 5) = N(1, 4) + 2 \cdot 3^{15 - 4} = N(1, 4) + 2 \cdot 3^{11}$. Note that by our definition of $N(n, k)$ to evaluate $N(1, 4)$ we just need to find the number of triples $a,b,c \leq 3^4$ such that $3 | a^3 + b^3 + c^3$. We can use Fermat's Little Theorem(which states that $a^3 \equiv a\mod 3$ so this would imply all residues that satisfy $a + b + c \equiv 0\mod 3$) or we can just list out all the residues. Ultimately, we get

 $a,b,c$ can be congruent to $1$ permutation of $0,0,0$ modulo $3$, $1$ permutation of $1,1,1$ and $1$ permutation of $2,2,2$. Then $6$ permutations of $0,1,2$ giving $1 + 1 + 1 + 6 = 9$ total permutations each having $\frac{3^4}{3} = 3^3$ possibilities. Thus we have

 $9 \cdot (3^3)^3 = 9 \cdot 3^9 = 3^{11}$. Therefore, $f(7, 6) = 3^{11} + 2 \cdot 3^{11} = 3 \cdot 3^{11}$. Now we have just found $f(7, 6)$ we still need to find $g(7, 6)$ which by our definition is

 $g(7, 6) = 2 \cdot 3^{18 - 7} = 2 \cdot 3^{11}$

 Finally, the answer is $N = 3 \cdot 3^{11} + 2 \cdot 3^{11} = 5 \cdot 3^{11}$ which gives an answer of $\boxed{735}$.

 




\end{enumerate}